\documentclass[14pt]{article} 
\usepackage[utf8]{inputenc}
\usepackage{epsfig}
\usepackage[ngerman]{babel}
\usepackage{mathtools,extarrows}
\usepackage[thmmarks]{ntheorem}
\usepackage{amsmath,theorem}
\usepackage{amsfonts}
\usepackage{amssymb}
\usepackage{txfonts}
\usepackage{marvosym}
\usepackage{MnSymbol}
\usepackage{mathrsfs}
\usepackage{parallel}
\usepackage{tikz}
\usepackage{geometry}
\usepackage{titletoc}
\usepackage{url}
\usepackage[pdftex,bookmarks=true,bookmarksnumbered=true,bookmarksopen=true,colorlinks=true,
		filecolor=black,linkcolor=blue,urlcolor=blue,plainpages=false,pdfpagelabels,
		citecolor=black,setpagesize=false]{hyperref}


\makeatletter
\renewcommand\l@section{\@dottedtocline{1}{0em}{2.5em}}
\renewcommand\l@subsection{\@dottedtocline{2}{2.5em}{2em}}
\makeatother

\titlecontents{chapter}[1.5em]{\addvspace{1em}\bfseries}
	{\contentslabel{1.5em}}{\hspace^{-1.5em}}
	{\hfill\contentspage}
\dottedcontents{section}[3.8em]{}{1.3em}{1pc}
\dottedcontents{subsection}[6.0em]{}{2.2em}{1pc}
\dottedcontents{subsubsection}[9.2em]{}{3.1em}{1pc}

%%%%theoremenvironments
\theoremstyle{plain} 

\theoremheaderfont{\bfseries}

\theorembodyfont{\it}
\newtheorem{theorem}{Theorem}[section]
\newtheorem{lem}[theorem]{Lemma}
\newtheorem{cor}[theorem]{Corollar}
\newtheorem{prop}[theorem]{Proposition}

\theorembodyfont{\rm}
%\newtheorem{remark}[theorem]{Bemerkung}
\newtheorem{defi}[theorem]{Definition}
\newtheorem{rul}[theorem]{Regel}

\theoremstyle{nonumberplain}

\theoremheaderfont{\bfseries}
\theorembodyfont{\normalfont}
\newtheorem{remark}{Bemerkung:}

\theoremheaderfont{%
    \normalfont\scshape}
\theorembodyfont{\normalfont}
\theoremsymbol{\square}
\newtheorem{proof}{Beweis:}

\topmargin=-2mm \textwidth=155mm \textheight=230mm
\oddsidemargin=0mm \evensidemargin=0mm
\thispagestyle{empty}
\newcommand{\mb}{\mathbb}
\newcommand{\mc}{\mathcal}
\newcommand{\mf}{\mathfrak}
\newcommand{\mr}{\mathrm}
\newcommand{\ms}{\mathscr}
\newcommand{\K}{\mb{K}}
\newcommand{\C}{\mb{C}}
\newcommand{\R}{\mb{R}}
\newcommand{\Q}{\mb{Q}}
\newcommand{\Z}{\mb{Z}}
\newcommand{\N}{\mb{N}}
\newcommand{\ra}{\rightarrow}
\newcommand{\LRa}{\Leftrightarrow}
\newcommand{\Ra}{\Rightarrow}
\newcommand{\La}{\Leftarrow}
\newcommand{\eps}{\varepsilon}
\newcommand{\vt}{\vartheta}
\newcommand{\vp}{\varphi}
\newcommand{\qq}{\hspace{.25cm}}
\newcommand{\ol}{\overline}
\newcommand{\ul}{\underline}
\renewcommand{\d}{ \, \mr{d}}
\newcommand{\Res}{\mr{Res}}
\newcommand{\Const}{\mr{Const} \, }
\newcommand{\Id}{\mr{Id}}
\renewcommand{\star}{\filledstar}
\newcommand{\Gl}{\mr{Gl}}
\newcommand{\sign}{\mr{sign}\,}
\newcommand{\dist}{\mr{dist}}
\newcommand{\diag}{\mr{diag}}
\newcommand{\CP}{\mr{CP}}
\newcommand{\CV}{\mr{CV}}
\newcommand{\supp}{\mr{supp}\,}
\newcommand{\RV}{\mr{RV}}
\renewcommand{\div}{\mr{div}\,}
\newcommand{\conv}{\mr{conv}}
\newcommand{\diam}{\mr{diam}}
\newcommand{\CS}{\mr{CS}}
\geometry{hmargin=3cm,top=2cm,bottom=4cm}

\newcommand{\makelicense}{Dieser Mitschrieb steht unter der freien \texttt{CC-BY-SA-DE 3.0} 
Lizenz.
\begin{center}\includegraphics[height=31pt,width=88pt]{license.pdf}\end{center}
Für weitere Informationen besuchen Sie
\begin{center}\url{http://creativecommons.org/licenses/by-sa/3.0/deed.de}\end{center}}


\input xy
\xyoption{all}

\title{\sc{Nichtlineare Funktionalanalysis}}
\author{Luis Felipe Müller}
%\date{Sommersemester 2011}

\begin{document}


\maketitle

\hrule \hspace*{1cm}\\[3mm]
{\sc Kompletter Mitschrieb zur gleichnamigen Vorlesung bei Herrn Dondl\\
\hspace*{1cm}\hfill (Sommersemester 2011, Uni Heidelberg)}\\
\hrule
\vspace{1cm}
\noindent\makelicense

%\vfill
%\listoffigures
%\vfill
%\listoftables

\newpage

\pagenumbering{roman}

\section*{Organisatorisches}

\subsection*{Termine}
\begin{itemize}
    \item Vorlesung: Mo./Mi. 9-11ct. -104 Ang. Math.
    \item Übungsaufgaben: Mi. - Mi. vor der Vorlesung. Kasten in der Angew.Math
    \item Übungsgruppe: Fr. 16-18 Angew. Math. -101
    \item Website zur Vorlesung: \url{http://dondl.org/wiki/Sommersemester_11}
    \item Literatur: \begin{enumerate}
        \item Růžička M: Nichtlineare Funktionalanalysis, Eine Einführung
        \item Aubin-Ekeland: Applied nonlinear Analysis
        \item Deimling: Nonlinear Functional Analysis
        \item Schwartz: Nonlinear Functional Analysis
        \item Zeidler: Nonlinear Functional Analysis and its applications 
        \end{enumerate}
    \item Prüfungen: Zulassung mit etwa 50\% der Übungsaufgaben-Punkte. Prüfung ist mündlich, 
            beispielsweise am 29. Juli (Fr.) 
    \item Dozent: Patrick Dondl, Sprechstunde Mo./Mi. 11-12 in Raum 130 (Angew. Math.)
    \item Tutor: Julian Scheuer
\end{itemize}

\newpage

\tableofcontents

\newpage

\pagenumbering{arabic}

\chapter{Einleitung}

\section{Thema der Vorlesung}

In der linearen Funktionalanalysis haben wir eine Vielzahl von Methoden kennengelernt um
Ergebnisse aus der endlichdimensionalen linearen Algebra auf den unendlichdimensionalen Fall 
zu verallgemeinern. Ein Hauptaufgabe war dabei, die Lösbarkeit von Gleichungen der Form
\[
    Ax=y
\]
für lineare Operatoren $A$ auf $\infty$-dimensionalen Banchräumen zu zeigen.
\begin{description}
    \item{Ein Beispiel:}
    
    Sei $\Omega\subset\R^3$, gefüllt mit einer inkompressiblen, riskosen Flüssigkeit.
    $v_j(x)$ sei die Geschwindigkeit der Flüssigkeit an der Stelle $x\in \Omega$. 
    $p(x)$ ist der Druck an der Stelle $x$.
    \begin{description}
    \item{Randbedingungen:}
    $v_j(x)=0 \qq x\in \partial\Omega$
    \item{Inkompressibilität:}
    $\partial_jv_j(x)=0 \qq x\in \Omega$
    \item{Bewegungsungleichung:}
    Wir betrachten die Kräfte, die auf einen kleinen Würfel, eingeschlossen durch
    $(x_1,x_2,x_3)$, $(x_1+\Delta x_1, x_2+\Delta x_2, x_3+\Delta x_3)$. Druck auf eine Oberfläche
    des Würfels mit Normale $x$, ist $f_j^p=p\cdot \Delta x_2 \Delta x_3 \cdot \delta_{ij} $ auf die
    gegenüberliegende Seite wirkt
    \[
        -(p+\partial _i p \Delta x_i) \Delta x_2 \Delta x_3 \delta_{ij}
    \]
    \end{description}
Zusammen ergibt sich 
\[
    f=(\partial_j p)\Delta V
\]

Kraft durch Viskosität auf eine Oberfläche mit Normale $x_1$ ist 
\[
    f_j^{V,x_i}= -2\eta \Delta x_2\Delta x_3 \partial_1v_j
\]
mit einer Konstante $\eta$. Der gleiche Trick wie oben ergibt für die gegenüberliegende
Oberfläche
\[
    \eta \Delta x_2 \Delta x_3 \partial (v_j+\partial v_j\Delta x_1)
\]
Zusammen ergibt sich
\begin{align*}
    f_j^V&= \eta\Delta V \cdot \partial_i\partial_iv_j\\
    \text{Newton:} \qq \rho \Delta V \frac{\d}{\d t} v_j (t,x(t))
        &= \eta \Delta V  \partial _i\partial_iv_j- \Delta V (\partial _j p) + \Delta 
        V\kappa_j \qq \text{mit einer externen Konstante $\kappa_j$, bspw. Gravitation}
\end{align*}
Teilen durch $\Delta V$ und die Kettenregel ergibt
\begin{align*}
    \rho \partial_t v_j&= \eta \partial_i\partial_i v_j - \rho (v_i\partial_i) - \partial_j p + K_j 
        \qq \text{(Newton)}\\
    d_jv_j&=0
\end{align*}
Naurier-Stokes. Frage: Existiert eine eindeutige Lösung zur sationären Naurier-Stokes-Gleichung:
\begin{align}\label{1}
    \eta \partial_i\partial_iv_j-(v_i\partial_i)v_j+\partial_jp+K_j=0
\end{align}
\end{description}
Wir können die Gleichung etwas umschreiben. Sei $H$ ein Hilbertraum
\[
    H:=\ol{\left\{ u\in C^\infty_c (\Omega,\R^3), \text{ so dass } \partial_jv_j=0 \right\}}
        ^{W^{1,2}(\Omega, \R^3)}
\]
Ein Skalarprodukt auf $H$ ist gegeben durch
\begin{align*}
    (u,v)_H:= \int_{\Omega} \nabla u\cdot \nabla v \d x = \int_\Omega \sum_{i,j=1}^3 
    (\partial_iu_j)(\partial_iv_j)
\end{align*}
$\Omega$ beschränkt $\Ra$ $(\cdot, \cdot)_H$ ist äquivalent zum üblichen Skalarprodukt 
(mittels Poincaré).
Wir multiplizieren (\ref{1}) mit $\omega \in H $, integrieren und erhalten
\begin{align*}
    \int_\Omega (\eta \partial_i\partial_i v_j- (v,\partial_i) v_j + K_j) \cdot \omega_j 
    &= \int_\Omega (\partial _jp)\omega_j=0 \text{, da $\omega$ divergenzfrei.}\\
    (\ref{1}) \Ra \eta (v,\omega)_H-a(v,v,w)-\int_\Omega K\omega&=0
\end{align*}
Ebenso für $a$:
\begin{align*}
    a(u,v,w):= (\underbrace{B(u,v)}_{\text{bilinear.}},w)_H
\end{align*}
Also
\begin{align*}
    (\ref{1}) \Ra \, (\eta v - B(u,v)) - \tilde K , \omega)_H=0 \qq \forall \omega \in H
\end{align*}
somit
\[
    \eta v - B(v,v) = \tilde K
\]
Das ist eine Gleichung der Form
\begin{align}\label{2}
    Fv=\tilde K,\qq\text{mit $F$ einem Nichtlinearen Operator}
\end{align}
Im ersten Teil der Vorlesung beschäftigen wir uns mit der eindeutigen Lösbarkeit von
Gleichungen der Form
\[
    Fx=y, \qq F:X\ra Y, \qq X,Y\qq \text{Banachräume}
\]
und zum Abschluß zeigen wir mit Hilfe des \textit{Schauder'schen Fixpunktsatzes} die Existenz
und finden eine Lösung von (\ref{2}), also der schwachen Form der stationären 
Navier-Stokes-Gleichung.

Im zweiten Teil der Vorlesung beschäftigen wir uns mit Variationsrechnung (d.h. dem Finden
von Minimierern nichtlinearer Funktionalen)
\[
    W:X\ra \R \qq\text{mit} \qq X \qq \text{ein Banachraum}
\]
Finde 
\[
    x_0\in X: \, W(x_0)= \inf_{y\in X}W(y)
\]
Insbesonder treffen wir dort auf Probleme in der Elastizitätstheorie.

\section*{Aufbau der Vorlesung}
\begin{itemize}
    \item Abbildungsgrad $\ra$ Existenz von Lösungen von $Fx=y$
    \item Monotone Operatoren $\ra$ Eindeutigkeit von Lösungen von $Fx=y$; 
    	zeitabhängige Probleme.
    \item Variationsrechnung $\ra$ $\inf_{y\in X} W(y)$
\end{itemize}

\subsection{Vorarbeiten}

\subsubsection{Ableitung in Banachräume und implizite Funktionen}

Es seien $X$ und $Y$ Banachräume, $\Omega\subset X$ offen, $F:\Omega \ra Y$, $x_0\in \Omega$

\begin{defi}[Gâteaux-Ableitung]
    Die Gâteaux-Ableitung $\d F(x_0,\psi)$ des Operators $F$ im Punkt $x_0$ in Richtung $\psi\in X$ 
    ist gegeben durch 
    \[
        \d F (x_0,\psi)=\lim_{s\ra 0} \frac{F(x_0+s\cdot \psi)-F(x_0)}{s}= \left. \frac{\d}{\d s} 
        F(x_0+s\psi)\right|_{s=0}
    \]
    falls der Limes existiert. Der Operator $F$ heißt in diesem Fall in $x_0$ Richtung $\psi$ 
    Gâteaux-differenzierbar.
\end{defi}
\begin{defi}[Fréchet-Ableitung]
    Der Operator $F$ heißt Fréchet-differenzierbar in $x_0\in \Omega$, falls ein beschränkter,
    linearer Operator 
    \[
        F'(x_0):X\ra Y 
    \]
    existiert, so dass 
    \begin{align}\label{3}
        \lim_{\|h\|\ra 0} \frac{\| F(x_0+h) - F(x_0) -F'(x_0)\cdot h \|  }{\|h\|} = 0
    \end{align}
    $F'(x_0)$ heißt dann Fréchet-Ableitung von $F$ in $x_0$.
\end{defi}

\begin{theorem}
    \begin{description}
        \item{i)} $F'(x_0)$ ist durch (\ref{3}) eindeutig bestimmt.
        \item{ii)} Falls $F$ stetig ist in $x_0$, so ist jeder lineare Operator, der (\ref{3}) erfüllt,
                ebenfalls stetig.
        \item{iii)} Ist $L:X\ra Y$ linear, so gilt
                \[
                    L'(x)=L \qq \forall x\in X
                \]
    \end{description}
\end{theorem}

\begin{proof}

\begin{description}

\item{i)} 
Es gelte (\ref{3}) auch für L. Dann haben wir
\begin{align*}
    \|Lh-F'(x_0)h\| \leq \eps \|h\| \qq \text{falls} \qq \|h\| < \delta = \delta(\eps, x_0)
\end{align*}
Für beliebiges $h$ folgt aber
\begin{align*}
    \|(L-F'(x_0)) (\delta\|h\|^{-1}\cdot h)\|&\leq \delta\eps\\
    \Ra \, \|(L-F'(x_0))h  \|&\leq \eps \|h\| \qq \forall \, h\in X, \, \forall \, \eps>0 \\
    \Ra \, \|L-F'(x_0)\|_{\ms L(X,Y)} &= 0
\end{align*}

\item{ii)}
(\ref{3}) wird umgeformt zu
\[
    \|F'(x_0) h \| \leq \eps \|h\| + \|F(x_0+h)-F(x_0)\|
\]
Mit $h\ra 0$ folgt die Stetigkeit (für $\|h\|\leq \delta$) von $F'(x_0)$ an der Stelle 0. 
Wegen Linearität von $F'(x_0)$ ist $F'(x_0)$ somit stetig. 

\item{iii)} (\ref{3}) gilt offensichtlich für $L'(x_0)=L$, mit i) folgt Eindeutigkeit.
\end{description}
\[ \]
\end{proof}


\begin{prop}
    Jeder Fréchet-differenzierbare Operator $F$ ist Gâteaux-differenziebar $\forall \, \psi\in X$ 
    und es gilt
    \[
        F'(x_0)\psi=\d F(x_0,\psi)
    \]
\end{prop}
\begin{proof}
    Übungsaufgabe \[\]
\end{proof}

\begin{defi}
    $F$ heißt (Fréchet-)differenzierbar auf $\Omega$, falls $\forall\, x\in X$ ein $F'(x)$ existiert,
    sodass $F'(x)$ stetig ist  und (\ref{3}) erfüllt. $F$ heißt stetig (Fréchet-)differenzierbar in
    $\Omega$, falls die Abbildung
    \[
        F': \Omega\ra \ms L(X,Y)
    \]
    stetig ist.
\end{defi}

\begin{prop}
    Existiert die Gâteaux-Ableitung $\d F(x,\psi) \, \forall x\in \Omega$, und ist sie linear und
    stetig in $\psi$ $\forall\, x\in \Omega$, so ist $F$ Fréchet-differenziebar auf $\Omega$ und
    es gilt 
    \[
        F'(x)\psi=\d F(x,\psi)
    \]
\end{prop}

\begin{proof} 
    Übungsaufgabe \[\]
\end{proof}

\begin{defi}
    Sei $F$ auf $\Omega$ stetig differenziebar, $x_0\in \Omega$. Falls ein stetiger linearer Operator
    \[
        F''(x_0):X\ra \ms L(X,Y)
    \]
    existiert mit
    \[
        \lim_{\|h\|\ra 0}\frac{\|F'(x_0+h)-F'(x_0)-F''(x_0)h\|_{\ms L(X,Y)}}{\|h\|}=0
    \]
    dann heißt $F$ in $x_0$ zweimal (Fréchet-)differenzierbar und $F''(x_0)$ heißt zweite Ableitung
    von $F$ in $x_0$. Höhere Ableitungen entsprechend.
\end{defi}

\begin{remark}
    Es gilt die Kettenregel:
    Seien $X$, $Y$, $Z$ Banachräume, $\Omega_X\subset X$ offen, $x_0\in \Omega_X$,
    \[
        F:\Omega_X\ra Y, \qq F(x_0)=y_0\in \Omega_Y\subset Y \qq \text{offen}
    \]
    \[
        G:\Omega_Y\ra Z
    \]
    Falls $F'(x_0)$ und $G'(y_0)$ existiert, so ist 
    \[
    (G\circ F)'(x_0)=G'(y_0)\circ F'(x_0)
    \]
\end{remark}

\begin{defi}[Partielle Ableitung]
    Seien $X,$ $Y$, $Z$ Banachräume, $\Omega_X\subset X$ offen, $x_0\in \Omega_X$, $\Omega_Y\subset Y$ 
    offen, $y_0\in \Omega_Y$. Der Operator
    \[
        F: \Omega_X \times \Omega_Y \ra Z 
    \]
    heißt partiell in $(x_0,y_0)$ nach dem zweiten Argument (nach $y$) differenzierbar, falls die 
    Ableitung 
    \[
        F(x_0, \cdot) : \Omega_Y \ra Z 
    \]
    differenziebar ist. Wir nennen den linearen Operator $F_Y(x_0,y_0): Y\ra Z$, der
    \[
        \lim_{\|h\|\ra 0} \frac{\| F(x_0,y_0+h) - F(x_0,y_0) - F_Y( x_0,y_0)h \|}{\|h\|} = 0 
    \]
    erfüllt, die partielle Ableitung von $F$ in $(x_0,y_0)$ nach dem zweiten Argument.
\end{defi}

\begin{prop}
    Seien $X,Y$ Banachräume, $\Omega \subset X $ offen und konvex mit $x_0, x_1\in \Omega$. 
    $F: \Omega \ra Y$ sei stetig Fréchet-differenzierbar auf $\Omega$. Dann gilt
    \[
        F(x_1)-F(x_0)=\int_{0}^1 F'(x_0+t(x_1-x_0))(x_1-x_0)\d t
    \]
    Das Integral ist als Limes der entsprechenden Riemannsumme zu verstehen und dieser existiert.
\end{prop}

\begin{proof}
    Übungsaufgabe \[\]
\end{proof}

Ähnlich dem endlichdimensionalen Fall geben uns die Ableitungen im Banachraum
hinreichende Bedingungen um Operatoren implizit zu definieren. Die Fragestellung ist die folgende:
Seien $X,\, Y,\, Z$ Banachräume, $U$ eine Umgebung von $x_0\in X$. $V$ eine Umgebung von $y_0\in Y$.
Wir suchen zu $F:U\times V \ra Z$ einen Operator
\[
    T:U_0\subset U \ra V, 
\]
sodass gilt
\[
    F(x,Tx)=F(x_0,y_0) \qq \forall \, x\in U_0.
\]
Durch eine einfache Verschiebung ist es ausreichend, den Fall
\[
    F(x_0,y_0)=0
\]
zu untersuchen.

\begin{prop}
    Sei $X$ ein Banachraum, $\Id: X\ra X$, $x\mapsto x$ die Identität auf $X$. Es sei
    \[
        R:B_r(0)\subset X \ra X, 
    \]
    eine $k$-Kontraktion mit $k<1$, d.h. $\|R(x)-R(y)\|\leq k \|x-y\|$, und es gelte
    \[
        \|R(0)\| < r(1-k)
    \]
    Dann existiert genau ein $x\in B_r(0)$ mit
    \[
        (\Id + R)x = 0
    \]
\end{prop}

\begin{proof}
    Sei $S=-R$, wir suchen also einen Fixpunkt von $S$.
    \begin{description}
        \item{1. Eindeutigkeit:} 
        Seien $Sx=x$ und $Sx'=x'$, damit gilt
        \[
            \|x-x'\|=\|Sx-Sx'\| \leq k \|x-x'\|.
        \]
        Mit $k<1$ folgt $x=x'$.
        
        \item{2. Existenz:}
        Sei $x\in B_r(0)$, es gilt
        \[
            \|Sx\|\leq \|Sx-S(0)\|+\|S(0)\|\leq k\|x\|+\|S(0)\|< kr + r(1-k) = r
        \]
        Sei $x_p=Sx_{p-1}$, $x_0=0$. Es gilt (wie auch im Banach'schen Fixpunktsatz, siehe ÜB 1), dass
        \[
            \|x_{n+p}-x_n\| \leq k^n(1-k)^{-1}\|x\|,
        \]
        damit ist $(x_p)_{p\in \N}$ eine Cauchy-Folge und konvergiert gegen $x\in X$.
        Wir haben weiter, dass
        \[
            \|x\| \leq \underbrace{\|x-x_{p+1}\|}_{\ra 0} + \|x_{p+1}\| \qq \text{mit} 
            \qq \|x_{p+1}\|< (1-k) \|S(0)\|=r\\
            \Ra \, \|x\|<r
        \]
        Wegen $x_{p+1}=Sx_p$ gilt dass $x=Sx$, somit ist $x$ der gesuchte Fixpunkt. 
    \end{description}
    \[  \]
\end{proof}

\begin{theorem}[Satz über implizite Funktion]\label{1.12}
    Seien $X$, $Y$, $Z$ Banachräume, $U\subset X$ Umgebung von $x_0\in X$, $V\in Y$ Umgebung von 
    $y_0 \in Y$. Sei weiter
    \[F:U\times V\ra Z\]
    stetig und stetig differenzierbar nach der zweiten Variablen. $F_Y(x_0,y_0)$ sei eine Bijektion
    von $Y$ nach $Z$ und es gelte
    \[
        F(x_0,y_0)=0
    \]
    Dann existiert $B_\delta(x_0)\subset U$, $B_r(y_0)\subset V$ und genau ein Operator 
    $T:B_\delta(x_0)\ra B_r(y_0)$, so dass $T(x_0)=y_0$ und $F(x,Tx)=0$ $\forall \, x\in B_\delta(x_0)$.
    $T$ ist stetig.
\end{theorem}

\begin{proof}
    Ohne Einschränkung sei $x_0=y_0=0$. Sei $L:= F_Y(0,0)$, $\Id: Y\ra Y$ die Identität auf $Y$. Es 
    sei $S(x,y):= L^{-1}F(x,y)-y$ für $(x,y)\in U\times V$. Somit gilt
    \[
        F(x,y)=0 \qq \LRa \qq y+S(x,y)=0.
    \]
    $S$ ist stetig differenzierbar nach dem zweiten Argument mit
    \[
        S_Y=L^{-1} F_Y(x,y)-\Id
    \]
    Damit gilt
    \[
        S_Y(0,0)=0
    \]
    Sei $k\in (0,1)$. Wegen Stetigkeit von $S_Y$ existiert $r>0$ mit 
    \[
        \|S_Y(x,y)\|\leq k \qq \forall (x,y)\in B_r(0)\times B_r(0)
    \]
    Sei nun $x\in B_r(0)$, $y, \tilde y \in B_r(0)$
    Es gilt nach Proposition \ref{1.10}, dass
    \[
        \|S(x,y)-S(x,\tilde y)\|=\left\|\int_0^1 S_Y(x,\tilde y +t(y-\tilde y)) (y-\tilde y) \d t 
        \right\|\leq k\cdot\|y-\tilde y\|
    \]
    Wegen $S(0,0)$ und Stetigkeit von $S$ existiert $\delta \leq r$, so dass
    \[
        \|S(x,0)\|\leq r(1-k) \qq \forall \, x \in B_\delta (0)
    \]
    Sei also $x\in B_\delta(0)$. Nach Proposition \ref{1.11} existiert genau ein $y\in B_r(0)$ mit
    $y+S(x,y)=0$. Wir setzen
    \[
        Tx=y, \qq T:B_\delta(0)\ra B_r(0)
    \]
    Es gilt $T(0)=0$ wegen $0+S(0,0)=T(0)+S(0,T(0))=0$ und der Eindeutigkeit von $T$. Es bleibt die
    Stetigkeit von $T$ zu zeigen: Seien $x,x'\in B_\delta (0)$, damit gilt
    \[
        0=Tx+S(x,Tx)=T(x'+S(x,Tx'))
    \]
    also
    \begin{align*}
        \|Tx-Tx'\|&\leq \|S(x',Tx')-S(x,Tx')\|+\|S(x,Tx)-S(x,Tx)\|\\
            &\leq \|S(x',Tx') - S(x,Tx')\|+k\|Tx-Tx'\|\\
            &= (1-k)\|Tx-Tx'\|\\
            &\leq \|S(x',Tx')-S(x,Tx')\| \ra 0 \qq \text{für $x\ra x'$} 
    \end{align*}
    Somit ist $T$ stetig.
    \[\]
\end{proof}

\begin{remark}
    Ist $F$ $r$-mal stetig differenziebrar, so gilt das auch für $T$.
\end{remark}

\begin{proof}
    Übungsaufgabe \[ \]
\end{proof}

\begin{theorem}\label{1.13}
    Seien $X$, $Y$ Banachräume, $U\subset X$ eine Umgebung von $x_0$. Es sei $F:U\ra Y$ stetig 
    differenzierbar und $F'(x_0)$ sei eine lineare Bijektion von $X$ nach $Y$.
    Dann existiert eine Umgebung $U_0\subset U$ von $x_0$, so dass
    \[
        F|_{U_0}: U_0\ra F(U_0)\ni y_0 -F(x_0)
    \]
    ein Homöomorphismus (bistetige Abbildung) ist.
\end{theorem}

\begin{proof}
    Wir wenden Satz $\ref{1.12}$ auf
    \[
        \tilde F(x,y):= F(x)-y
    \]
    an.
    \[  \]
\end{proof}

\begin{remark}
    Ist $F$ $r$-mal stetig differenzierbar, so gilt das auch für $F^{-1}$ ($F$ ist ein 
    $r$-Diffeomorphismus).
\end{remark}

\begin{proof}
    Übungsaufgabe\[ \]
\end{proof}

\begin{defi}[Zusammenhände Mengen]\label{1.14}
    \begin{description}
    \item{-}
    Sei $X$ ein (topologischer metrischer, normierter) Raum. Eine Menge $\Omega \subset X$ heißt
    zusammenhängend, falls es keine zwei abgeschlossenen (offenen) $\Omega_1$, $\Omega_2$ gibt
    mit
    \[
        \Omega\subset\Omega_1\cup\Omega_2, \qq \Omega\cap \Omega_1 \cap \Omega_2 = \varnothing, \qq
        \Omega\cap\Omega_{1,2}\neq \varnothing
    \]
    \item{-}
    Eine Menge $\Omega\subset X$ heißt wegzusammenhängend, falls sich je zwei Punkte in $\Omega$ durch
    eine stetige, in $\Omega$ verlaufende Kurve verbinden lassen.
    \item{-}
    Eine Menge $\ol\Omega \subset \Omega$ heißt Zusammenhangskomponente von $\Omega$, falls
    $\ol \Omega \subset \Omega$ maximal, zusammenhängend.
    \end{description}
\end{defi}

\begin{remark}
    Wegzusammenhängend $\Ra$ Zusammenhängend.

    \noindent Offen, zusammenhängend $\Ra$ Wegzusamenhängend
\end{remark}

\begin{theorem}[Mittelwert]\label{1.15}
    Seien $X$, $Y$ Banachräume, $F:X\ra Y$ stetig differenzierbar.
    \begin{description}
    \item{i)}
    Falls $\Omega$ konvex ist, so gilt
    \[
        \|F(x)-F(y)\|\leq M\|x-y\|,
    \]
    wobei
    \[
        M=\max_{0\leq t \leq 1}\| F'((1-t)x + ty) \|
    \]
    \item{ii)}
    Umgekehrt gilt: Falls
    \[
        \| F(x)-F(y) \|\leq M\|x-y\| \qq \forall \, x,y\in \Omega
    \]
    Dann gilt
    \[
        \sup_{x\in \Omega}\|F'(x)\|\leq M
    \]
    \end{description}
\end{theorem}

\begin{proof}
    Sei $f(t):= F((1-t)x+ty)$, $0\leq t\leq1$. Nach Kettenregel gilt
    \[
        f'(t)=F'((1-t)x+ty)(x-y)  
    \]
    \[
        \Ra \, \|f'(x)\|\leq\tilde M:= M\|x-y\|  
    \]
    \begin{description}
    \item{i)}
    Sei $\phi(t):=\|f(t)\|$ für $\delta>0$. Wir wollen zeigen, dass $\phi(t)\leq0$ $\forall \, 
    \delta>0$, $0\leq t\leq1$. Sei also (zum Widerspruch)
    \[
        t_0:=\max\{t\in [0,1]\, | \, \phi(s)\leq 0 \, \forall \, s\leq t\}.
    \]
    Dann gilt
    \begin{align*}
        \phi(t_0+\eps)&=\|f(t_0+\eps)-f(t_0)+f(t_0)-f(0)\|-(\tilde M + \delta)t\\
        &\leq \|f(t_0+\eps)-f(t_0)\|-(\tilde M +\delta) -\phi(t_0)\\
        &\leq \|f'(t_0)\eps+ \mc o (1)\|-(\tilde M +\delta)\eps\\
        &\leq (-\delta + \mc o (1))\eps 
    \end{align*}
    \item{ii)}
    Angenommen, es existiert $x_0$ mit $\|F'(x_0)\|\geq M+2\delta$, $\delta>0$. Dann existiert
    $e\in X$, $\|e\|=1$, $\| F'(x_0)e \|\geq M+\delta$. Somit gilt
    \begin{align*}
        M\eps&\geq \|F(x_0+\eps e) - F(x_0)\|= \|F'(x_0)(\eps e)+\mc o (\eps)\|\\
         &\geq (M+\delta)\eps- \mc o(\eps) > M\eps
    \end{align*}
    Das ist ein Widerspruch.
    \end{description}
    \[\]
\end{proof}

\begin{cor}\label{1.16}
    Sei $\Omega \subset X$ offen, (weg-)zusammenhängend, $F$ stetig differenzierbar auf $\Omega$.
    Es gilt
    \[
        F= \Const \qq \LRa \qq F'=0
    \]
\end{cor}
\begin{remark}
    Wir schreiben wie im endl. dim. $C(\Omega)= C^0(\Omega), \, C^1(\Omega)\, …$
\end{remark}

\noindent \textbf{Anwendungen}: Lokale Existenz und Eindeutigkeit Banachraum-wertiger
Differenztialgleichunegen. Sei $X$ Banachraum, $\Omega\subset X$ offen, $I\subset \R$
kompaktes Intervall. Es sei $C_b(I,\Omega)$ der Banachraum der beschränkten, stetigen Abbildungen
von $I$ nach $\Omega$, versehen mit der $\sup$-Norm.

\begin{lem}
    Sei $Y$ ein Banachraum, $f\in C(\Omega,Y)$ und sei die Funktion
    \[
        f_\star: C_b(I,\Omega)\ra C_b(I,Y)
    \]
    definiert als
    \[
        (f_\star x)(t)= f(x(t))
    \]
    Es gilt $f_\star\in C^r$
\end{lem}

\begin{proof}
    \begin{description}
    \item{$r=0$:}
    Sei $x_0\in C_b(I,\Omega)$, $\eps>0$ $\forall \, t \in I$ existiert $\delta(t)>0$, so dass
    \[
        \|f(\chi)-f(x_0(t))\|\leq \frac{\eps}{2} \qq \forall \, \chi\in U
    \]
    Die offenen Kugeln 
    \[
        \{ B_{\delta(t)}(x_0(t)) \}_{t\in I}
    \]
    sind eine offene Überdeckung vom $\{ x_0(t) \}_{t\in I}$. Diese Menge ist als stetiges Bild einer
    kompakten Menge kompakt, und somit existiert endliche Teilüberdeckung
    \[
        \{ B_{\delta(t_j)}(x_0)(t_j) \}_{1\leq j \leq N}
    \]
    Sei nun $x\in C_b(I,\Omega)$ mit
    \[
        \|x-x_0\|\leq \delta := \min_{1\leq j \leq N}\delta(t_j)
    \]
    Somit existiert $\forall \, t\in I$ ein $t_j$, so dass $\|x_0(t)-x_0(t_j)\|\leq \delta(t_j)$,
    und deshalb gilt
    \[
        \|f(x(t))- f(x_0(t))\|\leq \|\underbrace{f(x(t))-f(x_0(t_j))}_{\leq 2\delta}\|
        +\|\underbrace{f(x_0(t_j))-f(x_0(t))}_{\leq\delta}\|,
    \]
    denn
    \[
        \|x(t)-x_0(t_j)\|\leq \|x(t)-x_0(t)\|+\|x_0(t)-x_0(t_j)\|\leq 2\delta(t_j).
    \]
    Somit folgt die Steigkeit……
    \item{$r=1$:}
    Wir müssen zeigen, dass
    \[
        \sup_{t\in I} \|f(x_0(t))+x(t)- f(x_0(t)) - f'(x_0(t))x(t)\|\leq \eps \sup_{t\in I}\|x(t)\|
    \]
    denn
    \[
        (f'_\star(x_0)x)(t)= f'(x_0(t))x(t)
    \]
    Übungsaufgabe. Folgt wie Stetigkeit durch Kompaktheit von $I$.
    \end{description}
    \[  \]
\end{proof}

\section{Der Brouwer'sche Abbildungsgrad}

\subsection*{Motivation}
\begin{description}
    \item[Ziel:]
    $f(x)=0$ zu lösen für $f:U\subset X \ra X$, $X$ Banachraum.
    \item[Frage:]
    Existenz/Anzahl der Lösungen
    \item[Rückblick auf Funktionentheorie:]
    Sei $z_0\in \C$
    \[
        n(\gamma,z_0) = \frac{1}{2\pi i} \int_\gamma \frac1{z-z_0} \d z
    \]
\end{description}
$\leadsto$ Verallg.: $f\in \mc H(\C)$. $0\nin f(\gamma)$
\[
    n(f(\gamma),0)=\frac1{2\pi i} \int_\gamma \frac{f'}{f} \d z = \sum_k n(\gamma,z_k) \alpha_k
\]
wobei $f(z_n)=0$, $\alpha_k$ Vielfachheiten.

\noindent Ziel: Verallg. des Begriffs „Umlaufzahl“ für Abb. $f:U\subset \R^n\ra \R^n$

\subsection{Die Determinantenformel}

\subsubsection{Notation}

$U\Subset \R^n$ offen, $f\in C^1(U,\R^n)$
\begin{align*}
    J_f(x)&:=\det \d f(x)\\
    \RV(f)&:=\{ y\in \R^n \,| \, \forall \, x\in f^{-1}(y), \, J_f(x)\neq 0 \}\\
    \CV(f)&:= \R^n \setminus \RV(f)\\
    D_y^r(\ol U, \R^n)&:=\{f\in C^k(\ol U, \R^n)\,|\,y\nin f(\partial U)\}\\
    D_y(\ol U, \R^n)&:=D_y^0(\ol U, \R^n)
\end{align*}
$\tau(\R^n)$ bezeichne die Topologie auf $\R^n$.
\begin{defi}
    Eine Abbildung
    \[
        \deg: \bigcup_{U\in \tau(\R^n),\, y\in \R^n} (D_y(\ol U, \R^n)\times \{U\}\times \{y\}\ra \R,
    \]
    d.h.
    \[
        \deg=\deg(f,U,y) 
    \]
    heißt Gradabbildung, falls
    \begin{description}
        \item[(D1)]
        Translationsinvarianz: 
        \[
            \deg(f,U,y)=\deg(f-y,U,0)
        \]
        \item[(D2)]
        Normalisation:
        \[
            \deg(\Id,U,y)=1 \qq \forall y\in U
        \]
        \item[(D3)]
        Additivität: Seien $U_1,U_2\subset U$ offen und disjunkt, sodass
        \[
            y\nin f(\ol U\setminus(U_1\cup U_2)),
        \]
        dann gelte
        \[
            \deg(f,U,y)=\deg(f,U_1,y)+\deg(f,U_2,y)
        \]
        \item[(D4)]
        Homotopieinvarianz: 
        \[
            H(t)=(1-t)f+tg \in D_y(\ol U,\R^n) \qq \forall \, t\in [0,1]\, \Ra \, \deg(f,U,y)=\deg(g,U,y)
        \]
\end{description}
\end{defi}

\begin{theorem}\label{2.1}
    Sei $\deg$ eine Gradabbildung. Dann gilt
    \begin{description}
        \item{(i)}
        $\deg(f,\varnothing,y)=0$ und
        \[
            \deg(f,U,y)=\sum_{i=1}^N \deg(f,U_i,y)
        \]
        falls 
        \[
            y\nin f(\ol U\setminus\bigcup_{i=1}^NU_i),
        \]
        wobei $U_i\subset U$ offen und disjunkt für $1\leq i \leq N$.
        \item{(ii)}
        $y\nin f(U) \, \Ra \, \deg(f,U,y)=0$
        \item{(iii)}
        $|f(x)-g(x)|<\dist(y,f(\partial U)) \, \forall \, x\in \partial U \, \Ra \, \deg(f,U,y)
        =\deg(g,U,y)$
    \end{description}
\end{theorem}

\begin{proof}
    \begin{description}
    \item{(i)}
    Sei $U_1=U$, $U_2=\varnothing$, einsetzen in \textbf{(D3)}
    \[
        \Ra \, \deg(f,\varnothing,y)=0
    \]
    $i=1$: $U_2=\varnothing$
    \[
        \Ra \, \deg(f,U,y)=\deg(f,U_1,y)
    \]
    $i>1:$ Induktion mittels \textbf{(D3)}
    \item{(ii)}
    \begin{align*}
        y\nin f(U) \Ra y &\nin f(\ol U)\, \Ra \, y\nin f(\ol U\setminus \varnothing)\\
        \overset{(i)}{\Ra } \, \deg (f,U,y)&=0 \qq (i=1, U_1=\varnothing)
    \end{align*}
    \item{(iii)}
    Sei $H(t,x):=(1-t)f(x)+tg(x)$ und sei $x\in \partial U$
    \begin{align*}
        \Ra \, |H(t,x)-y|&= |f(x)-y+t(g(x)-f(x))|\\
            &\geq |f(x)-y| - |g(x)-f(x)|\\
            &\geq \dist(y,f(\partial U))- |g(x)-f(x)| >0\\
        \Ra \, y&\nin H(t,\partial U) \qq \forall \, t\,\Ra \,H(t)\in D_y(\ol U,\R^n)\\
        &\overset{(\textbf{D4})}{\Ra} \text{ Behauptung.}
    \end{align*}
    \end{description}
    \[ \]
\end{proof}

\begin{theorem}\label{2.2}
    \begin{description}
        \item{i)} $\deg(\cdot,U,y)$ ist lokal konstant in $D_y(\ol U,\R^n)$
        \item{ii)} $\deg(f,U,\cdot)$ ist lokal konstant in $\R^n\setminus f(\partial U)$
        \item{iii)} Seien $H:[0,1]\times \ol U\ra \R^n$ und $y:[0,1]\ra \R^n$ stetig (d.h. $H$ ist eine
            Homotopie zwischen $H(0)=H(0,\cdot)$ und $H(1)$), so gilt
        \[
            \deg(H(0),U,y(0))=\deg(H(1),U,y(1)),
        \]
        falls $H(t)\in D_{y(t)}(\ol U, \R^n)$ $\forall t \in [0,1]$
    \end{description}
\end{theorem}

\begin{proof}
    Beachte: $D_y(\ol U, \R^n)$ ist offen in $C^0(\ol U, \R^n)$
    \begin{description}
    \item{i)}
        $\|f-g\|_{C^0,\ol U} <\eps \, \Ra \, |f(x)-g(x)|<\eps \, \forall \, x\in\partial U$ mit
        \[
            \eps:= \dist(y,f(\partial U)) \, \Ra \, \deg (f,U,y)=\deg (g,U,y)
        \]
    \item{ii)}
        Sei $y_0\nin f(\partial U)$ und $y\in B_{\dist(y_0,f(\partial U))}(y_0)\subset \R^n\setminus
                f(\partial U)$
        \begin{align*}
            \Ra \, \|(f-y)-f\|&<\dist(y_0,f(\partial U))\\
                \overset{i)}{\Ra} \, \deg(f-y,U,y_0) &= \deg(f,U,y_0)\\
            \overset{\textbf{(D1)}}{\Ra } \, \deg (f,U,y_0+y)&= \deg(f,U,y_0)
        \end{align*}
    \item{iii)}
        $H$ ist gleichmäßig stetig
        \[
            \Ra \, H:[0,1]\ra C^0(\ol U,\R^n)\qq t\mapsto H(t,\cdot)
        \]
        ist auch stetig. $H$ ist ein stetiger Weg in $D_y(\ol U,\R^n)$. Sei $y$ fest
        \[
            \Ra \, \deg(H(t,\cdot),U,y)=\Const,
        \]
        weil $\deg(\cdot,U,y)$ konst. auf Zusammenhangskomponenten ist. Für $y=y(t):$
        \begin{align*}
            \deg(H(0),U,y(0))&=\deg(H(0)-y(0),U,0)=\deg(H(t)-y(t), U, 0) \qq \forall t\in [0,1]\\
            &= \deg(H(1),U,y(1)).
        \end{align*}
    \end{description}
    \[ \]
\end{proof}

\begin{lem}\label{1.3.5}
    Zwei nichtsinguläre Matrizen $M_1,M_2\in \Gl(n)$ sind genau dann homotop in $\Gl(n)$, falls
    \[
        \sign \det M_1 = \sign\det M_2
    \]
\end{lem}

\begin{proof}
    \begin{description}
    \item{„$\Ra$“} Sei $M\in \Gl(n)$. Wegen der Linearität von $\det $ in Zeilen können elementare
    Zeilenumformungen mit Hilfe stetiger Deformationen in $\Gl(n)$ erzeugt werden.
    \[
        M \, \leadsto\diag(m_1,…,m_m), \qq \text{mit} \qq |m_i|=1
    \]
    weil $\sign \det M_1=\sign \det M_2$.
    \[
        H(t):= \begin{pmatrix} \pm \cos (\pi t)& \mp \sin(\pi t)\\ \sin(\pi t) & \cos(\pi t)\end{pmatrix}
    \]
    ist eine Homotopie in $\Gl(n)$ von $\diag(\pm 1,1)$ nach $\diag(\mp 1,-1)$. $i=n$ transformiere
    $\diag(\mp 1,-1)$.

    \noindent $i=n$ transformiere $\diag(m_i,…,m_i)$ nach $\diag(\pm1,1)$

    \[
        \leadsto \begin{pmatrix} \sign\det M& 0& \cdots&0\\ 0& 1 &\cdots&0\\\vdots&&\ddots&0\\ 0&\cdots&
        \cdots& 1\end{pmatrix}
    \]

    \end{description}
\end{proof}

\begin{theorem}
    Sei $f\in D_y^1(\ol U, \R^n)$, $y\nin \CV(f)$ und $\deg$ eine Gradabbildung. Dann gilt
    \[
        \deg(f,U,y)=\sum_{x\in f^{-1}(y)} \sign J_f(x),
    \]
    wobei die Summe endlich ist.
\end{theorem}

\begin{proof}
    O.B.d.A. $y=0$ (nach \textbf{D1}). Alle $x\in f^{-1} (0)$ sind isolierte Punkte in $U$
    (Homöomorphiesatz).
    $f^{-1}(0)$ hat höchstens am Rand einen Häufungspunkt, aber $0\nin f(\partial U)$.
    \[
        \Ra \, f^{-1}(0)= \{ x^i \}_{i=1}^N
    \]
    Wähle $\delta >0$ so klein, dass $B_\delta(x^i)$ paarw. disjunkt.
    \[
        \overset{\text{\ref{2.1} i)}}{\Ra} \, \deg(f,U,0)=\sum_{i=1}^N \deg(f,B_\delta(x^i),0).
    \]
    Beachte $0\nin f(\ol U \setminus \bigcup_{i=1}^N B_\delta(x^i))$.
    \begin{align*}
        f\in C^1 \, \Ra \, f(x)&= \d f(x)(x-x^i) + |x-x^i| r(x-x^i) \qq \text{mit} \qq r\in C^0(B_\delta(x^i),\R^n), 
        \qq r(0)=0.\\
           H(t,x) &:= \d f (x^i) (x-x^i)+ (1-t) (x-x^i)+(x-x^i)
    \end{align*}
    Zeige $0\nin H(t,\partial B_\delta(x^i))$, für alle $t\in [0,1]$.
    \[
        J_f(x^i)\neq 0 \, \Ra \, \exists\, \lambda >0:\, |\d f(x^i)(x-x^i)|\geq \lambda|x-x^i|
    \]
    O.B.d.A. sei $\delta$ so klein, dass $|r(x-x^i)|<\lambda$ in $B_\delta (x^i)$.
    \begin{align*}
        \Ra\, | H(t,x) |&>|\d f(x^i)(x-x^i)|-(1-t)(x-x^i)r(x-x^i)\geq \lambda \delta - \delta|r|>0
        \qq \forall \, x\in \partial B_\delta (x^i)\\
        \overset{(\textbf{D4})}{\Ra} \, \deg(f,U,0)&=\sum_{i=1}^N\deg (\d f (x^i)(\cdot-x^i),
                B_\delta(x^i),0)
    \end{align*}
    Lemma \ref{1.3.5}
    \[
        \Ra \, \deg(\d f (x^i)(\cdot- x^i),B_\delta (x^i),0)= \deg(\diag(\sign J_f(x^i),1,…,1),
            B_\delta(x^i),0)
    \]
    Falls $J_f(x^i)>0\, \overset{\textbf{(D2)}}{\Ra}\,\deg(I(\cdot-x^i),B_\delta (x^i),0)=1$. Es genüngt
    also, $\deg(M(\cdot-x^i),B_1(x^i),0)$. Zu berechnen, wobei $M=\diag(-1,1,…,1)$ ist.
    \begin{align*}
        U_1&:= B_1 (x^i) = \{ \max_{1\leq k\leq n} |x_k-x_k^i|<1 \}\\
        U_2&:= U_1+ (2,0,\cdot,0)\\
        g(r)&=2-|r-1|, \, h(r)=1-r^2\\
        f_1(x)&:= (1-g(x_1-x_1^i)h(x_2-x_2^(i))…h(x_n-x_n^{(i)}),…, 1)\\
        f_2(x)&:= (1,x_2-x_2^{(i)}, … , x_n-x_n^{(i)})\\
        f_1^{-1}(0)&= \{ y,z \} \qq y=x^i, z = x^i+(2,0,…,0)\\
        f_1|_{\partial U}&= f_2|_{\partial U} \qq \Ra \qq \deg(f_1,U_2,0)=0\\
        \Ra \, \deg (f_1,U,0)&= \deg(f_1,U_1,0)+ \deg(f_2,U_2,0) \qq (\star)\\
        \Ra \, \deg(M,B_1(x^i),0)&=\deg(\partial f_1(y),B_1(x^i), 0)\\
        &=\deg(f_1, U_1,0) \overset{(\star)}{=} -\deg (f_1,U_2,0)\\
        &= -\deg(\partial f_1(z),U_2,0)=\deg(\Id, U_2,0)=-1
    \end{align*}
    \[ \]
\end{proof}

\subsection{Verallgemeinerung der Determinantenformel}

Wir haben einen Kandidaten für $\deg$ identifiziert; für $f\in D'_Y(\ol\Omega,\R^n)$, $y\in \R^n\cap
RV(f)$ gilt
\[
    \deg(f,\Omega, g)=\sum_{x_j\in f^{-1}(y)} \sign J_f(x_j).
\]

\noindent \textbf{Problem:} Nur für glatte $f$, und reguläre Punkte zu definieren. Programm zum 
Existenzbeweis: Verallgemeinerung der Determinanten-Formel auf
\begin{description}
    \item{-} kirtische Werte $y$.
    \item{-} nur stetige Funktionen $f$.
\end{description}

\noindent Drei Beispiele zur Illustration:
\begin{enumerate}
    \item $f(x)=x^2$, $U=(-1,1)$
    \item $f(x)=x^2$, $U=(-1,2)$
    \item $f(x)=x+2\sin(x)$, $U=(-10,10)$
\end{enumerate}
Wie aus dem Beispielen ersichtlich ist, haben wir noch (kleine) Schwierigkeiten, den Abbildungsgrad
an kritischen Werten von $f$ zu definieren. Stetige Fortsetzungen liegt aber nahe. Das funktioniert
aber nur falls es von diesen nicht „zu viele“ gibt.

\subsubsection*{Schritt 1:} Kritische Werte von $f$:

\begin{lem}[Sard]\label{2.5}

Es sei $f\in C^1(\Omega, \R^n)$, $\Omega\subset\R^n$ ist offen und beschränkt. Dann ist
$CV(f)$ eine Lebesgue-Nullmenge.
\end{lem}
    
\begin{proof}
    Klar, falls $f$ eine affine Abbildung ist. (Dimension des Bildraumes der (konstanten) Ableitung!)
    Wir linearisieren. 

    \noindent Es sei $\CP(f):=\{x\in \Omega\,|\, J_f(x)=0\}$ die Menge der kritischen Punkte von $f$.
    Sei $\{ Q_j \}_{j\in \N}$ eine abzählbare, offen Überdeckung von $\Omega$ bestehend aus Würfeln,
    so dass
    \[
        \ol \Omega_i\subset \Omega, \, i\in \N.
    \]
    Es gilt 
    \[
        CV(f)=f(\CP(f))=\bigcup_{j\in \N} f(\CP(f)\cap Q_j)
    \]
    Es reicht also zu zeigen, dass
    \[
    |f(\CP(f)\cap Q_j)|
    \]
    für alle $j\in \N$ verschwindet. Sei nun $Q$ ein solcher Würfel, $\rho$ dessen Kantenlänge.
    Sei $\eps >0$, und sei $Q$ unterteilt in $N^n$ Würfel $Q^i$ der Seitenlänge $\frac{\rho}{N}$, so
    dass
    \begin{align}\label{4}
        |f(x)-f(\tilde x)-f'(\tilde x)(x-\tilde x)|&\leq \int_0^1|f'(\tilde x+t(x-\tilde x))-
        f'(\tilde x)|\cdot|x-\tilde x|\d x\leq \frac{\eps\rho}{N} \qq \text{für alle} 
            \qq x,\tilde x \in Q^i
    \end{align}
    So ein $N$ existiert, nachdem $f'$ auf $Q$ gleichmäßig stetig ist. Nun enthalte $Q^i$ einen
    kirtischen Punkt $\tilde x_i \in \CP(f)$, ohne Einschränkung sei $\tilde x_i=0$, $f(\tilde x_i)=0$,
    und $M=f'(\tilde x _i)$. Wegen $\det M =0$ existiert eine ONB $\{b^j\}_{j=1}^n$ mit
    \[
    b^n \bot R(M)
    \]
    Weiter gilt 
    \begin{align*}
        Q^i&\subset \left\{ \sum_{j=1}^n \lambda _j b^j \, \big| \, \| \lambda \|_2\leq \sqrt{n} 
            \frac{\rho}{N} \right\}\subset\left\{\sum_{j=1}^n\lambda _jb^j \, \big|
            \, |\lambda_j|\leq \sqrt{n}\frac{\rho}{M} \, \forall 1\leq j\leq n\right\}
    \end{align*}
    Damit existiert $C>0$, (unabhängig von $i$), so dass
    \[
        MQ^i\subset \left\{\sum_{j=1}^{n-1}\lambda_jb^j\, \big| \, |\lambda_j|
            \leq C\cdot \frac{\rho}{N} \right\}
    \]
    mit $C=\sqrt{n} \max_{x\in \ol Q}|f'(x)| $. Damit gilt nach (\ref{4}) sogar, dass
    \[
        f(Q^i)\subset\left\{ \sum_{j=1}^n \lambda_jb^j\, \big| \, |\lambda_j|\leq 
        (C+\eps)\frac{\rho}{N} \, \forall 1\leq j\leq n-1, \, |\lambda_n|\leq \frac{\eps\rho}{N}\right\}
    \]
    Es folgt
    \[
        |f(Q^i)|\leq \frac{\tilde C\eps}{N^n},
    \]
    falls in $Q^i$ ein kritischer Punkt liegt. Es gibt maximal $N^n$ Unterwürfel $Q^i$ mit kritischen
    Punkten. Somit gilt
    \[
        |f(Q\cap \CP(f))|\leq C\cdot \eps
    \]

\end{proof}

\noindent Dank Lemma \ref{2.5} ist $\R^n\setminus \CV(f)$ dicht in $\R^n$. Das reicht leider (?)
noch nicht.

\noindent Wir brauchen $d_1=d_2$, um den Abbildungsgrad sinnvoll durch die Determinanten-Formel
definieren zu können. (denn $\deg$ soll konstant sein auf Zusammenhangskomponenten, unabhängig
von krit. Werten.)

\textbf{Idee:} Umschreiben der Determinantenformel als \textit{Integral}.
Es sei im Weiteren $\eta_\eps$ ein Standard-Mollifier (Standard-Dirarcfolge), d. h.
\[
    \eta_\eps \in C^\infty_c(B_\eps(0)\subset \R^n), \qq \int_{\R^n} \eta_\eps=1,\qq \eta_\eps\geq 0
\]

\begin{lem}\label{2.6}
    Sei $f\in D^1_y(\ol \Omega, \R^n)$, $y\nin \CV(F)$. Dann gilt
    \begin{align}\label{5}
        \deg(f,\Omega,y)&=\sum_{x_j\in f^{-1}(y)} \sign J_f(x_j)\\
                        &= \int_\Omega \eta_\eps (f(x)-y) J_f(x)\d x
    \end{align}
    für alle $\eps$ hinreichend klein, d. h.
    \[
        \eps_0>\eps>0, \qq \text{mit} \qq \eps_0=\eps_0(f,y)
    \]
    Es gilt $\supp (\eta_\eps(f(\cdot)-y))\subset \Omega$ für $\eps < \dist(y,f(\partial \Omega))$
\end{lem}

\begin{proof}
    \begin{description}
        \item{1)} 
        Falls $f^{-1}(y)=\varnothing$, dann sei $\eps_0=\dist(y,f(\partial \Omega))$
        \item{2)}
        Falls $f^{-1}(y):= \{x^i\}_{1\leq i \leq N}$, sei $\eps_0>0$, so dass
        \[
            f^{-1}(B_{\eps_0}(y)) = \bigcup_{i=1}^N U_i
        \]
        mit $U_i\cap U_j=\varnothing$ für $i\neq j$.
        Aus dem \textit{Satz über die imliziete Funktion} (war ja klar, dass wir den mal brauchen)
        folgt für evtl. noch kleiners $\eps_0$, dass
        \[
            f\big|_{U_i(x_i)} \qq \text{bijektiv,} \qq J_f(x)\neq 0 \qq \forall x \in U_i(x_i).
        \]
        Wieder gilt $\eta _{\eps}(f(x)-y)=0$ für
        \[
            x\in \ol \Omega \setminus \bigcap _{i=1}^N U_i(x^i)
        \]
        Damit gilt
        \begin{align*}
            \int_\Omega \eta_\eps(f(x)-y) J_f(x) \d x&= \sum_{i=1}^N \sign J_f(x^i)\cdot
            \int_{B_{\eps_0(x)}}\eta_\eps (\tilde x ) \d \tilde x\\
            &=\sum_{x^i\in f^{-1(y)}} \sign  J_f(x^i)
        \end{align*}
        mit $f(x)=\tilde x$.
    \end{description}
    \[ \]
\end{proof}
Die Integraldarstellung ergibt auch für kritische Punkte Sinn. Aber:
Wegen Abhängigkeit von $\eps_0$ von $f,y$ ist. Auch der Wert des Integrals nicht a priori stetig in 
$f,y$, denn $\eps_0$ hängt ab von $f$ und $y$.


Dieses Problem können wir beseitigen, wenn wir $f\in C^2$ fördern. Wir brauchen zunächst einen Hilfssatz.

\begin{prop}\label{2.7}
    \begin{description}
        \item{a)}
        Sei $u\in C^1(\R^n)$, $\supp u \subset K$, K kompakt. Sei $\vp=\div u$. Dann gilt
        \[
            \int_{\R^n} \vp (x) \d x=0
        \]
        \item{b)}
        Sei $\vp \in C^1(\R^n)$, und sei $z\in \R^n$. Dann gilt
        \[
            \vp (x+z)-\vp(x)=\div \left( z\cdot \int_0^1 \vp (z+tz) \d t\right)
        \]
        \item{c)}
        Seien $u$, $K$, $\vp$ wie in $a)$, $\Omega \subset \R^n$ offen, beschränkt.
        Sei $f\in C^2(\ol \Omega)$ mit $K\cap f(\partial \Omega)=\varnothing$. Dann existiert
        $v\in C^1(\ol \Omega)$, mit $\supp v \subset \Omega$, so dass
        \[
            \vp (f(x))J_f(x)=\div v(x) \qq \text{auf} \qq \Omega
        \]
    \end{description}
\end{prop}

\begin{proof}
    \begin{description}
    \item{a)}
    Klar!
    \item{b)}
    Sei
    \[
        \eta (x):= \int_0^1 \vp (x-sz) \d s.
    \]
    \begin{align*}
        (z\cdot \eta(x))&=\frac{\d}{\d t } \eta (x+tz)\big|_{t=0}=\int_0^1 \frac{\d}{\d t} \vp(x+sz+tz)
       \big|_{t=0} \d s\\
        &= \int_0^1 \frac{\d}{\d t} \vp(x+sz)\d s=\vp(x+z)-\vp(x)
    \end{align*}
    \item{c)}
    Es sei $\d _{ij}$ der $(i,j)$- Eintrag der Kofaktormatrix von $(f')_{ij}= \partial _if_j(x)$.
    Es sei $v_j(x)=\sum_{j=1}^n u_j(f(x))\d _{ij}(x),i=1,…,n$. Es gilt $K\cap f(\partial \Omega) = \Phi$,
    $f\in C(\ol \Omega)$, damit existiert $\delta>0$ mit $\dist(K,f(\ol \Omega \setminus 
                f(\Omega_\delta)))>0$.
    Somit gilt $\supp v\subset \Omega_\delta\subset \Omega$.
    Wir rechnen:
    \begin{align*}
        \partial_iv_i&= \sum_{\d _{ij}(x)\partial_ku_j(f(x))}\partial _i f_k.+\sum_{j=1}^nu_j(f(x))
        \partial_i \d_{ij}x
    \end{align*}
    \textit{Behauptung 1:}
    \[
        \sum_{i=1}^n \partial_i\d_{ij}(x)=0
    \]
    \textit{Behauptung 2:}
    \[
        \sum_{i=1}^n \d_{ij} (x) \partial_if_k(x) =\delta_{ij}\cdot J_f(x)
    \]
    Somit gilt:
    \begin{align*}
        \div u(x)&=\sum_{k,j}\partial _k u_j(f(x))\cdot \left( \sum_{i=1}^n \d_{ij}(x) \partial f_k(x) 
                \right) +\sum_{i=1}^nu_j(f(x)) \left( \sum \partial _i \d_{ij}(x) \right)\\
        =&\sum_{k,j}\partial _k u_j(f(x))\delta _{jk} J_f(x)=\vp(f(x))J_f(x)
    \end{align*}
    Behauptungen: Siehe Übungsblatt 2.
    Wir wollen nun zeigen, dass der über die Determinantenformel definierte Abbildungsgrad konstant
    ist auf Zusammenhangskomponente. Zu zeigen: $\deg(f,\Omega,y_1)=\deg(f,\Omega,y_2)$.

    \begin{lem}\label{2.8}
        Sei $f\in C^2(\ol \Omega,\R^n),$ $ y_0\nin f(\partial \Omega) $,
        \[
            \rho:=\dist(y_0,f(\partial \Omega))
        \]
        Dann ist $\deg(f,\Omega,\cdot)$ (Definiert durch die Determinantenformel) konstant auf
        \[
            B_\delta (y_0)\cap \RV(f)
        \]
    \end{lem}

    \begin{proof}
        Sei $y^j\in B_\delta (y_0)\cap \RV(f),$ $j=1,2$, sei $\delta := \rho-\max_{j=1,2}\|y^j-y_0\|$.
        Sei $\eps>0$, so dass
        \[
            \deg(f,\Omega,y^j)=\int \eta_\eps(f(x)-y^j)J_f(x) \d x \qq \text{(nach Lemma $\ref{2.6}$)}
        \]
        mit Proposition $\ref{2.7}$ b) gilt
        \begin{align*}
            \eta_\eps(x-y^2)-\eta_\eps(x-y^1)&= \eta_\eps (x-y^1+(y^1-y^2))-\eta_\eps(x-y^1)\\
            &= \div w(x) \qq\text{mit}\\
            w(x)&= (y^1-y^2)\int_0^1 \eta_\eps (x-y^1+t(y^1-y^2))\d t
        \end{align*}
    \textit{Behauptung:} Es gilt $\supp w\subset B_\delta(y_0)$
    \begin{proof}
        Sei $x\in \supp w$. Damit existert $t\in [0,1]$ mit $|x-y^1+t(y^1-y^2)|<\eps$
        \begin{align*}
            \Ra \, |x-y_0|&\leq \eps + |y^1-t(y^1-y^2)-y_0|\leq \eps +|(1-t)(y^1-y_0)+t(y^2-y_0)|\\
                \eps + \rho-\delta&<\rho
        \end{align*}
        Damit gilt aber, dass $f(\partial \Omega)\cap \supp w=\varnothing$. Mit Proposition 
        $\ref{2.7}$ e) existiert
        \[
            v\in C^1(\Omega),\, \supp v \subset \Omega.
        \]
        und
        \[
            \left( \eta_\eps(f(x)-y^2)-\eta_\eps (f(x)-y^1) \right)J_f(x)=\div v(x)
        \]
        Mit der Proposition $\ref{2.7}$ a) folgt die Behauptung. \[  \]
    \end{proof}

    \end{proof}
    \end{description}
\end{proof}

\begin{defi}\label{2.9}
    Sei $f\in C^2(\ol \Omega,\R^n)$, $y\in f(\partial \Omega)$. Wir setzen
    \[
        \deg (f,\Omega,y)= \deg(f,\Omega,\tilde y)=\sum_{x\in f^{-1}(\tilde y)} \sign J_f(x),
    \]
    wobei $\tilde y \in \RV(f)$ mit
    \[
        |\tilde y- y| < \dist(y,f(\partial \Omega))
    \]
\end{defi}

\noindent\textit{Behauptung:} $\deg$ ist somit nach dem vorherstehenden Überlegungen wohldefiniert.
\subsubsection*{2. Schritt: Nur stetige Funktionen $f$}

Auch hier die Idee: Sei $f\in D_y(\ol \Omega, \R^n)$. Wir suchen $\tilde f \in C^2 (\ol \Omega, \R^n)$
hinreichend nahe an $f$ (in derselben Zusammenhangskomponente von $D_y(\ol \Omega, \R^n)$) und übertragen
den Wert von $\deg(\cdot , \Omega, y)$ von $\tilde f$ auf $f$. Es bleibt zu zeigen, dass die so
definierte Funktion stetig ist.

\begin{lem}\label{2.10}
    Sei $f\in D_y^2(\ol \Omega, \R^n)$, sei $g\in C^2(\ol \Omega, \R^n)$. Dann existiert $\eps >0$, so 
    dass
    \[
        \deg (f+tg,\Omega,y)= \deg(f,\Omega,y) \qq \forall t\in (-\eps,\eps)
    \]
\end{lem}
\begin{proof}
    \begin{description}
        \item{1)}
        $f(y)=\varnothing \, \Ra \, (f+tg)^{-1}(y)=\varnothing$ falls
        \[
            |t| < \frac{\dist(y,f(\ol \Omega))}{\|g\|_\infty}
        \]
        \item{2)}
        $y\in \RV (f) \, \Ra \, f^{-1}(y)=\{x\}_{i=1}^N$. Mithilfe des \textit{Satzes über die 
            implizite Funktion} finden wir
            \[
                U(x^i)=:U^i
            \]
            disjunkte Umegbungen, so dass eindeutige Lösungen $x^i(t)\in U^i$ existieren von
            \[
                (f+tg)(x)=y \qq \forall |t|<\eps
            \]
            Wir können (zumindest auf evtl. noch kleineren $U^i$) annehmen, dass das Vorzeichen
            von $J_{f+tg}$ konstant ist auf $U^i$. Sei 
            \[
                \eps _2= \frac{\dist(y,f(\Omega\setminus \bigcup_{i=1}^NU^i))}{\|g\|_\infty}.
            \]
            Dann gilt $\dist(y,(f+tg)(\partial \Omega))>0$ $\forall |t|<\eps_2$. Das Lemma gilt somit
            für $\eps = \min(\eps,\eps_2)$
        \item{3)}
            $y\in \CV(f)$. Sei dann $\tilde y\in \RV(f)$.
            \[
                |y-\tilde y|< \frac{\rho}{3} = \frac13 \dist(y,f(\partial \Omega)).
            \]
            Nach Definition \ref{2.9} gilt
            \[
                \deg (f,\Omega,y)=\deg(f,\Omega,\tilde y).
            \]
            Sei $\tilde \eps>0$ mit $\deg(f,\Omega,\tilde y)=\deg(f+tg,\Omega,\tilde y)$ für $|t|<\tilde
            \eps$ (Nach Schritt 2).
            Mit $\eps=\min(\tilde \eps , \frac{\rho}{3\|g\|_\infty})$ gilt:
            \[
                |\tilde y-(f+tg)(x)|\geq \frac{\rho}{3} \qq \forall x \in \partial \Omega
            \]
            Somit gilt
            \[
                |\tilde y-y|< \dist (\tilde y, (f+tg)(\partial \Omega))
            \]
            also
            \[
                \deg(f+tg,\Omega,\tilde y)= \deg (f+tg,\Omega,y)
            \]
            Es folgt
            \[
                \deg(f,\Omega, y)= \deg(f+tg,\Omega,y)
            \]
    \end{description}
    \[ \]
\end{proof}

\begin{theorem}[Existenz und Eindeutigkeit des Brouwer'schen Abbildungsgrades]\label{2.11}
    Es existiert eine eindeutige Abbildung
    \[\deg(f,\Omega,y)\]
    mit den Eigenschaften \textbf{(D1)}-\textbf{(D4)}. Weiter gilt
    \[
        \deg(\cdot, \Omega, y): D_y(\ol \Omega, \R^n)\ra \Z
    \]
    ist konstant auf Zusammenhangskomponente von $D_y(\ol \Omega, \R^n)$. Für $f\in D_y(\ol \Omega,\R^n)$
    ist $\deg(f,\Omega,y)$ gegeben durch
    \[
        \deg(f,\Omega,y)=\sum_{x\in \tilde f ^{-1}(y)} \sign J_{\tilde f}(x),  
    \]
    wobei
    \[
        \tilde f \in D_y^2(\ol \Omega, \R^n)
    \]
    beliebig aus derselben Komponente von $D_y(\ol \Omega,\R^n)$ wie $f$ gewählt werden kann mit $y
    \in \RV(\tilde f)$.
\end{theorem}

\begin{proof}

Aus Lemma \ref{2.8} und \ref{2.10} folgt, dass $\deg$ wohldefiniert ist und lokal konstant mit 
Werten in $\Z$. Die Abbildung $\deg$ auf Zusammenhangskomponenten.
\begin{description}
\item[(D2)] $\ra$ klar.
\item[(D1)] gilt nachdem diese Bedingung per Konstruktion für die Determinatenformel gilt.
\item[(D3)] Wir wählen $\|f-\tilde f\|_\infty<\dist(y,f(\ol \Omega\setminus(\Omega_1\cap \Omega_2)))$.
    \textbf{D3} gilt per Konstruktion für $\tilde f$, sonst für $f$.
\item[(D4)] folgt aus der Konstruktion von $f$ auf Zusammenhangskomponente.
\end{description}
\[ \]
\end{proof}

\subsubsection*{Beispiele:}
\begin{description}
    \item{1)}
    Sei 
    \[
        f(x_1,x_2)=\begin{pmatrix} x_1-2x_2+\cos(x_1+x_2)\\x_2+2x_1+\sin(x_1+x_2) \end{pmatrix}
    \]
    und
    \[
        g(x_1,x_2):=\begin{pmatrix} x_1-2x_2 \\ x_2+2x_1 \end{pmatrix}
    \]
    Es gilt $|g(x)|=\sqrt{5} |x|$ und $|f(x)-g(x)|=\sqrt{\sin^2(x_1+x_2)+\cos^2(x_1+x_2)}=1$.
    Sei $h(t)=(1-t)g+tf=g+t(f-g)$.
    \begin{align*}
        |h(t)|&\geq |g|-t|f-g|>0 \qq \text{für} \qq |x|>\frac1{\sqrt{5}}\\
        \Ra \, \deg(f,B_r(0),0)&= \deg(g,B_r(0),0)=1 \qq \text{für} \qq r>\frac1{\sqrt{5}}
    \end{align*}
    Somit existiert eine Lösung $x:f(x)=0$.

    \item{2)}
    \begin{theorem}
        Ein stetiges Vektorfeld im $\R^n$, das auf einer Kugeloberfläche überall nach außen zeigt, muss
        auf einem Punkt im Innern der Kugel verschwinden. 
        
        Anders formuliert, sei $f:\ol{B_R(0)} \ra \R^n$ stetig, so dass $f(x)\cdot x>0$ $\forall |x|=R$. 
        Dann existiert ein $x_0\in B_R(0)$ mit $f(x_0)=0$.
    \end{theorem}

    \begin{proof}
        Wir haben $\deg(\Id, B_R(0),y)=1$ für $y\in B_R(0)$. Angenommen, $f(x)\neq 0$ für alle
        $x\in B_R(0)$. Dann gilt $f^{-1}(0)\cap B_R(0)=\varnothing$.
        \[
            \Ra \, \deg(f,B_R(0),0)=0
        \]
        Sei $H(t)=(1-t)\Id+tf$.
        \begin{align*}
            \deg(H(0),B_R(0),0)&=1\\
            \deg(H(1),B_R(0),0)&=0.
        \end{align*}
        Es existiert ein $t_0\in (0,1)$, so dass $H(t_0)\nin D_0(\ol{B_R(0)},\R^n)$.
        \[
            \exists x_0\in \partial B_R(0)\cdot (H(t_0))(x_0)=0.
        \]
        \begin{align*}
            0&=H(t_0)(x_0)\\
            \Ra \, 0&= H(t_0)(x_0)\cdot x_0= \underbrace{(1-t_0)R^n}_{>0}+\underbrace{t_0f(x_0)x_0}_{>0}
            >0
        \end{align*}
        Widerspruch!
        \[ \]
    \end{proof}
\end{description}

\subsection{Der Brouwer'sche Fixpunktsatz}

Der \textit{Brouwer'sche Fixpunktsatz} ist eine Folgerung aus den Eigenschaften des Abbildungsgrades.
Er besagt, dass stetige Abblidungen, die kompakte, konvexe Mengen im $\R^n$ (oder Mengen, die dazu
homöomorph sind) in sich selbst abbilden, einen Fixpunkt besitzen.

\subsubsection{Einschub:}

\begin{description}
        \item[Zerlegung der Eins:]
        Sei $\{U_i\}_{i\in I}$ eine Überdeckung von $X$, $X$ ein topologischer Raum.
        $\{\vp_\lambda\}_{\lambda\in\Lambda}$ ist eine $\{U_i\}_{i\in I}$ untergeordnete lokal endliche
        Zerlegung der Eins, falls gilt
        \begin{itemize}
        \item[-] $\vp_\lambda \in C(U,[0,1]) \,\forall\, \lambda$.
        \item[-] $\forall \lambda\in \Lambda\, \exists \, i \in I: \, \supp \vp_\lambda\subset U_i$
        \item[-] $\sum_{\lambda \in \Lambda} \vp _\lambda=1$
        \item[-] $\forall x \in X \, \exists \, U(x)$ Umgebung von $x$, so dass nur
        \textit{endlich viele} $\lambda \in \Lambda$ existieren, mit $U(x)\cap \supp \vp_\lambda\neq 
        \varnothing$ (lokal endlich).
        \end{itemize}
        \item[Konstruktion] aus einer lokal endlichen Überdeckung $(V_\lambda)_{\lambda\in\Lambda}$ 
        ($\forall x\in X \, \exists \, U(x)$ Umgebung: $U(x)\cap V_\lambda\neq \varnothing$ nur für
         endlich viele $\lambda\in \Lambda$.) in metrischen Räumen $(X,d)$.
        \begin{itemize}
            \item[-] Sei $\alpha(x)=\sum_{\lambda\in \Lambda} \dist(x,X\setminus V_\lambda)>0 \, \forall
            x \in X$, mit $\dist(x,\varnothing):=1$.
            \item[-] $\vp _\lambda(x):=\frac{1}{\alpha(x)}\dist(x,X\setminus V_\lambda) \, \in [0,1]$,
            $\vp_\lambda=0$ für $x\nin V_\lambda$.
        \end{itemize}
        \item[„Konstruktion“] einer lokal endlichen Überdeckung $\{V_\lambda\}_{\lambda\in\Lambda}$, die
        eine offen Überdeckung $\{U_i\}_{i\in I}$ verfeinert (d.h. $\forall \, \lambda\in \Lambda \, 
        \exists i \in I, \, V_\lambda\subset U_i$).
        \item{Die Existenz} einer solchen lokal endlichen verfeinernde Überdeckung zu jeder offenen
        Überdeckung ist die Definition der \textit{Parakompaktheit}.
        \begin{theorem*}
                    Jeder metrische Raum ist parakompakt.
        \end{theorem*}

        \begin{idea}
        Angenommen $I=\N$ also wohlgeordnet. (jede Teilmenge besitzt ein eindeutiges kleinstes Element.)
        Wir setzen für $i\in I, \, n\in \N$.
        \[
            D_{in}=\bigcup_{x\in \Phi(i,n)} B_{2^{-n}}(x)
        \]
        \[
            \Phi(i,n):=\left\{x\in X \, \Big| \, \text{mit}\, \begin{cases}\text{$i$ ist die kleinste
            Zahl, so dass $x\in U_i$}\\ \text{$x\nin D_{jm}$ für $m<n$}\\ B_{3\cdot2^{-n}}\subset U_i  
            \end{cases}\right\}
        \]
        \begin{description}
        \item klar: $\{ D_{in} \}$ verfeinert $\{U_i\}$
        \item[Überdeckung:] Klar, denn für jedes $x\in X$ finden in ein kleinstes $i\in I=\N$ und ein
        $n$ hinreichend groß.
        \item[lokale Endlichkeit:] (s.Artikel).
        Ist $I$ nun eine beliebige Indexmenge, so läßt sich diese wohlordnen (\textit{Zorn'sches Lemma}),
        und die Wohlordnung ist die einzige Eigenschaft von $I$, die wir benutzt haben.
        \end{description}
        \[ \]
        \end{idea}
    \end{description}
\begin{theorem}[Fortsetzung stetiger Funktionen]\label{2.12}
    Sei $X$ ein metrischer Raum, $Y$ normierter Raum. Sei $K\subset X$ abgeschlossen. Sei $F\in C(K,Y)$.
    Dann besitzt $F$ eine stetige Fortsetzung $G:X\ra Y$, so dass
    \[
        G(X)\subset\conv\, F(K)
    \]
\end{theorem}

\begin{proof}
    Wir betrachten die offene Überdeckung
    \[
        \left\{ B_{\rho(x)}(x) \right\}_{x\in X\setminus K} \qq \text{von} \qq X\setminus K
    \]
    mit $\rho(x)=\frac12 \dist(x,K)$.
    Wir wählen nun eine lokal endliche Zerlegung der Eins $\{\vp_\lambda\}_{\lambda\in \Lambda}$, welche
    der offenen Überdeckung untergeordnet ist.
    Sei
    \[
    G(x):= \begin{cases} F(x)& \text{für} \qq x\in K\\
           \sum_{\lambda \in \Lambda } \vp_\lambda(x) \cdot F(x_\lambda) & \text{für} 
           \qq x\in X\setminus K\end{cases},
    \]
    wobei $x_\lambda$ beliebig in $K$ so gewählt ist, dass 
    \[
        \dist(x_\lambda,\supp \vp_\lambda)\leq 2\cdot \dist(K,\supp \vp_\lambda).
    \]

    $G$ ist offensichtlich stetig auf $X\setminus\partial K$ und
    \[
        G(X) \subset \conv\,G(K).
    \]
    Sei also $x_0\in\partial K$ und sei $\eps>0$. Wir wählen $\delta >0$, so dass
    \[
        \|F(x)-F(x_0)\|\leq\eps \qq \forall x\in K \qq \text{mit} \qq d(x,x_0)\leq 9\delta.
    \]
    Es bleibt zu zeigen, dass
    \[
        \|G(x)-F(x_0)\|\leq \eps \qq \forall x\in X\setminus K \qq \text{mit} \qq d(x,x_0)\leq \delta
    \]a
    Sei also $x\nin K$, dann gilt
    \[
        \|G(x)-F(x_0)\|\leq \sum_{\lambda\in \Lambda} \vp_\lambda (x) \|F(x_\lambda)-F(x_0)\|
    \]
    Nach Konstruktion liegt $x_\lambda$ nicht weit von $x$ entfernt, damit nicht weit von $x_0$, falls
    $x\in\supp \, \vp _\lambda$. In der Tat gilt für $x\in \supp\,\vp_\lambda$:
    \begin{align*}
        d(x,x_\lambda)&\leq \dist(x_\lambda,\supp\,\vp_\lambda)
        +\underbrace{\diam(\supp\, \vp_\lambda)}_{=\sup_{x,y\in\supp\,\vp_\lambda} d(x,y)}\\
            &\leq 2\dist(K,\supp \, \vp _\lambda)+\diam(\supp\, \vp _\lambda)
    \end{align*}
    Nachdem $\{\vp _\lambda\}_{\lambda\in\Lambda}$ der Überdeckung $\{B_{\rho(x)}\}_{x\in X\setminus K}$
    untergeordnet ist, existiert $\tilde x\in X\setminus K$, so dass
    \begin{align*}
        \supp\, \vp _\lambda &\subset B_{\rho(\tilde x)}(\tilde x)
    \end{align*}
    \begin{align*}
        \Ra \, \diam (\supp \, \vp _\lambda) &\leq 2\rho(\tilde x)=\dist(\tilde x,K)
        \leq2\dist(K,B_{\rho(\tilde x)}(\tilde x))\leq 2 \dist(K,\supp \vp_\lambda)\\
        \Ra \, d(x_0,x_\lambda)&\leq 4 \dist(K,\supp\, \vp_\lambda).
    \end{align*}
    Und es folgt
    \begin{align*}
        d(x_0,x_\lambda) &\leq d(x_0,x)+d(x,x_\lambda)\\
                &\leq d(x_0,x)+4\dist(K,\supp\,\vp_\lambda)\\
                &\leq d(x_0,x)+8\dist(K,\supp\, \vp_\lambda)\\
                &\leq d(x_0,x)+8d(x_0,x)=9d(x_0,x)\leq \delta.
    \end{align*}
    Nach Wahl von $\delta$ gilt
    \[
        \|F(x_{\lambda})-F(x_0)\|\leq \eps \qq \forall \lambda\in\Lambda:\, \vp_\lambda(x)\neq 0.
    \]
    Somit gilt
    \[
        \|G(x)-G(x_0)\|\leq \eps \qq \text{für} \qq d(x,x_0)\leq \delta.
    \]
\end{proof}

\begin{remark}
    Mit Hilfe dieses Satzes und dem Abbildungsgrad lässt sich der sog. „Igelsatz“ zeigen, der besagt,
    dass man einen Igel nicht stetig kämmen kann. Übungsaufgabe. 
\end{remark}

\begin{theorem}[Brouwerscher Fixpunktsatz]\label{2.13}
    Sei $K$ ein topologischer Raum, der zu einer kompakten konvexen Teilmenge des $R^n$ homöomorph ist.
    Sei $f\in C(K,K)$. Dann besitzt $f$ einen Fixpunkt.
\end{theorem}

\begin{proof}
    \begin{description}
    \item{1.}
    $K=\ol{B_r(0)}\subset\R^n$. Falls ein Fixpunkt auf dem Rand existiert, dann sind wir fertig.
    Ansonsten gilt für $H(t)=\Id-tf$, dass $0\nin H(t)(\partial B_r(0))$, nachdem
    \[
        |H(t)(x)|\geq |x|-t|f(x)|\geq (1-t)r>0 \qq \text{für} \qq 0\leq t< 1, \qq x\in\partial B_r(0).
    \]
    Nach der Annahme der Nichtexistenz eines Fixpunktes auf $\partial B_r(0)$ ist auch $H(1)(x)\neq 0$
    $\forall x \in \partial B_r(0)$.
    \[
        \Ra \, \deg(\Id-f,B_r(0),0)=\deg(\Id,B_r(0),0)=1.
    \]
    Somit existiert $x\in B_r(0)$ mit $x=f(x)$.
    \item{2.}
    Sei nun $K\subset \R^n$ konvex, kompakt. Für ein $\rho>0$ gilt $K\subset B_\rho(0)$ und gemäß Satz
    \label{2.12} können wir $f$ stetig durch $g$ auf $\ol{B_{\rho}(0)}$ fortsetzen mit
    \[
        g\left( \ol{B_\rho (0)} \right)\subset \conv \, K=K.
    \]
    Nach 1. finden wir $x\in \ol{B_\rho(x)}$ mit $x=g(x)$.
    Es gilt $g(x)\in K$, somit folgt $x\in K$ und wir haben $x \in K$ mit $f(x)=x$.
    \item{3.}
    Sei $K$ homöomorph zu $K^\star \subset \R^n$ kompakt, konvex. Sei $h$ die entsprechende Homöomorphie.
    nach 2. besitzt $h\circ f \circ h^{-1}$ einen Fixpunkt $x^\star\in K^\star$. Damit ist aber $x=h^{-1}
    (x^\star)\in K$ ein Fixpunkt von $f$.
    \end{description}
    \[ \]
\end{proof}

\begin{remark}
    \begin{description}
    \item{1.}
    Die Bedingungen sind tatsächlich notwendig. Gegenbeispiele siehe Übung.
    \item{2.}
    Es existiert auch eine stetige Abbildung $f\in C(\ol{B_1(0)},\ol{B_1(0)})$, $B_1\subset X$ separabler
    $\infty$-dimensionaler Hilbertraum ohne Fixpunkt.
    \begin{description}
    \item{$\ra$}
    Beispiel von Kakutani später.
    \item{$\ra$}
    \textit{Schauderscher Fixpunktsatz}.
    \end{description}
    \item{3.}
    Im eindimensionalen Fall ist \ref{2.13} nichts anderes als der Zwischenwertsatz angewendet auf
    $x-f(x)$.
    \item{4.}
    Vergleich mit dem \textit{Banach'schen Fixpunktsatz}: Wesentlich geringere Anforderungen an den
    Operator (nur Stetigkeit), dafür hohe Anforderungen an den Raum (endlichdim., kompakt, konvex).
    Wir bekommen keine Eindeutigkeitsaussage.
    \end{description}
\end{remark}

\subsubsection*{Ein Anwendugsbeispiel}

Existenz positiver (bzw. nichtnegativer) Eigenwerte und Eigenfunktionen. Sei $A=(a_{ij})_{i,j=1}^n$ eine
Matrix und sei $a_{ij}\geq 0$ $\forall (i,j)$. Dann existiert $\lambda\geq 0$, $x=(x_i)_{i=1}^n$,
$x_i\geq 0$ $\forall i$ mit
\begin{align}\label{6}
    Ax=\lambda x
\end{align}

\begin{proof}
    Sei 
    \[
        K=\{ x\in \R^n \, |\, x\geq 0, \, \sum_{i}x_i=1 \}
    \]
    kompakt und konvex.
    \begin{description}
        \item{1)}
        Falls $Ax=0$ für ein $x\in K$, gilt \ref{6} mit $\lambda=0$.
        \item{2)}
        Sei $Ax\neq 0$ $\forall x\in K$. Dann existiert $\alpha>0$ mit
        \[
            \sum_i(Ax)_i\geq \alpha \qq \forall x\in K.
        \]
        Es sei also
        \[
            f:x\mapsto \frac{Ax}{\sum_i (Ax)_i}; \qq (f(x))_i \geq 0 \qq \forall i
        \]
        und es gilt
        \[
            \sum_i (f(x))_i=1 \qq \forall x\in K.
        \]
        Dann ist $f(K)\subset K$ und nach Satz $\ref{2.13}$ existiert ein Fixpunkt $x\in K$ mit $x=f(x)$,
        d.h.
        \[
            Ax=\lambda x \qq \text{mit} \qq \lambda =\sum_i(Ax)_i >0
        \]
    \end{description}\[ \]
\end{proof}

\subsection{Weitere Eigenschaften des Abbildungsgrades}

Sei $m<n$, wir identifizieren im Folgenden den $\R^m$ mit dem Unterraum
\[
    \{x\in \R^n \, | \, x_{m+n}=…=x_n=0\}\subset \R^n.
\]

\begin{theorem}[Reduktionseigenschaften des Abbildungsgrades]\label{2.14}
    Sei $\Omega \subset \R^n$ offen und beschränkt, $\Omega \cap \R^m \neq \varnothing$ und sei
    \[
        f:\ol \Omega \ra \R^m\, \text{stetig,} \qq g=\Id-f
    \]
    Sei $y\in \R^m$, $y\nin g(\partial \Omega)$. Dann gilt
    \[
        \deg(g,\Omega,y)=\deg(g|_{\ol{\Omega}\cap\R^m},\Omega|_{R^m},y).
    \]
\end{theorem}

\begin{proof}
    Sei $\Omega_m=\Omega\cap\R^m\neq \varnothing$ offen, beschränkt im $\R^m$, und
    $g_m:=g|_{\ol{\Omega_m}}$. Es gilt
    \[
        \partial \Omega_m=\partial\Omega\cap\R^m.
    \]
    und
    \[
        g_m(\ol{\Omega_m})\subset\R^m,\, y\nin g_m(\partial\Omega_m).
    \]
    $\Ra$ $\deg(g_m,\Omega_m,y)$ ist definiert (als Abbildungsgrad im $\R^m$).
    \begin{description}
        \item{1)}
        Sei $g\in C^1(\ol \Omega)$ und $y\in \RV(g)$ und $x\in \Omega$, $g(x)=y$.
        \[
        \Ra \, x=y=f(x)\in \R^m \qq \Ra \qq x\in \Omega_m.
        \]
        Somit haben wir
        \[
            g^{-1}(y)=g_m^{-1}(y).
        \]
        Nach der Determinantenformel genögt es nun zu zeigen, dass
        \[
            J_{g_m}(x)=J_g(x) \qq \text{für} \qq x\in \Omega_m.
        \]
        Es sei $\Id_k$ die $k\times k$-Einheitsmatrix. Es gilt
        \begin{align*}
            J_{g_m}(x)&=\det(\Id_m-f'(x)) \qq (f|_{\Omega_m}')\\
            J_g(x)&= \det\begin{pmatrix} \Id_m-(\partial_jf_i(x)) & -\partial_jf_i\\
            0& \Id_{n-m}\end{pmatrix} \qq \text{wegen} \, f(\ol\Omega)\subset\R^m.
        \end{align*}
        Die gewünschte Aussage folgt sofort durch Entwicklung der Determinante nach den letzten
        $n-m$ Zeilen.
        \item{2)}
        Der allgemeine Fall folgt durch Wahl von $\tilde g=\Id-\tilde f$ hinreichend nah an $g$, so dass
        $\tilde g$ die in 1. geforderten Eigenschaften besitzt.
    \end{description}
    \[ \]
\end{proof}


Es sei nun wieder $\Omega \subset \R^n$ offen, beschränkt, $f\in C(\ol \Omega,\R^n)$. Nach Satz \ref{2.2}
gilt, dass $\deg(f,\Omega,\cdot)$ konstant ist auf Zusammenhangskomponenten von $\R^n\setminus f(\partial 
\Omega)$. Wir bezeichnen diese Zusammenhangskomponenten mit $\{G_j\}_{j\in I}$ und schreiben $\deg(f,
\Omega,y)=\deg(f,\Omega,G_j)$ für $y\in G_j$

\begin{theorem}[Produktregel]\label{2.15}
    Sei $\Omega\subset\R^n$ offen, beschränkt. Seien $G_j$ die Zusammenhangskomponenten von $\R^n
    \setminus f(\partial \Omega)$ für $f\in C(\ol\Omega, \R^n)$ und sei
    \[
        g\circ f \in D_y(\ol \Omega,\R^n), \qq y\in \R^n
    \]
    Dann gilt
    \begin{align}\label{9}
        \deg(g\circ f, \Omega, y)=\sum_j\deg(f,\Omega,G_j)\cdot \deg(g,G_j,y)
    \end{align}
\end{theorem}

\begin{proof}
    (Endlichkeit der Summe) Es gilt $f(\ol\Omega)$ ist kompakt, somit existiert $r>0$ mit $f(\ol \Omega)
    \subset B_r(0)$. Nachdem $g^{-1}(y)$ abgeschlossen ist, gilt $g^{-1}(y)\cap B_r(0)$ kompakt ist.
    $\{G_j\}_{j\in I}$ ist eine offene Überdeckung dieser Menge, es genügen somit endlich viele $\{  
    G_j \}_{j=1}^N$ um $g^{-1}(y)\cap B_r(0)$ zu überdecken.
    \begin{description}
    \item{1)}
    Wir nehmen an, dass $f\circ g\in C^1(\ol \Omega)$, $y\in \RV(g\circ f)$. Es gilt nach der
    Kettenregel, dass
    \[
        (g\circ f)'(x)=g'(f(x))\circ f'(x)
    \]
    Die Behauptung folgt dann durch eine längere Rechnung.
    \begin{align*}
        \deg(g\circ f,\Omega,y)&=\sum_{x\in(g\circ f)^{-1}(y)}\sign J_{g\circ f}(x)=\sum_{x\in (g\circ f)
        ^{-1} (y)}(\sign J_g(f(x)))\cdot(\sign J_f(x))\\
        &= \sum_{z\in g^{-1}(y)} \sign J_g(z)\cdot \sum_{x\in f^{-1}(z)} \sign J_f(x)
        =\sum_{z\in g^{-1}(y)} \sign J_g(z)\cdot \deg(f,\Omega,z)
    \end{align*}
    Mit der Überdeckung $\{G_j\}_{j\in I}$ von $g^{-1}(y)$ gilt
    \begin{align*}
        \deg(g\circ f,\Omega,y)&=\sum_{j=1}^m\sum_{z\in g^{-1}(y)\cap G_j}\sign J_gi(z)\cdot \deg(f,\Omega,z)
        =\sum_{j=1}^N \deg(f,\Omega,G_j) \cdot \sum_{z\in g^{-1}(y)\cap G_j} \sign J_g(z)\\
        &= \sum_{j=1}^N \deg(f,\Omega,G_j)\cdot \deg(g,G_j,y)
    \end{align*}
    Diese Formel gilt nun natürlich auch für $y\in \CV(g\circ f)$ und für $g$ nur stetig. Etwas
    problematischer ist der Fall, dass $f$ nur stetig ist, da sich bei der Modifikation von $f$ die
    Mengen $G_j$ ändern.

    Es sei
    \[
        L_l=\{z\in \R^n\setminus(\partial \Omega)\, | \, \deg(f,\Omega,z)=l\}
    \]
    Für $l\neq 0$ gilt, dass $L_l$ aus einer Vereinigung von Mengen $G_j$ bestehen muss. Sei nun
    $\tilde f\in C^1 (\ol \Omega)$, so dass $|f(x)-\tilde f(x)|< \frac 12 \dist(g^{-1}(y),
    f(\partial\Omega))$ für $x\in \ol \Omega$. Wir definieren $\tilde G_j$ und $\tilde L_l$ entsprechend
    für die Funktion $\tilde f$.
    Es gilt
    \[
        L_l\cap g^{-1}(y)=\tilde L_l\cap g^{-1}(y).
    \]
    nach Satz \ref{2.1} (iii). ($\|f(x)-\tilde f(x)\|<\dist(y,f(\partial \Omega)\qq \forall
    x\in \partial \Omega)\qq \Ra \qq \deg(f,\Omega,y)=\deg(\tilde f,\Omega,y)$.) Es gilt somit
    \begin{align*}
        \deg(g\circ f, \Omega, y)&=\deg(g\circ\tilde f,\Omega,y)\\
        &=\sum_{j}\deg(f,\Omega,\tilde G_j)\cdot \deg(g,\tilde G_j,y)=\sum_{l\neq0} l\cdot\deg( g,\tilde
        L_l,y)=\sum_{l\neq 0} l\cdot \deg(g,L_l,y)\\
        &=\sum_j \deg(f,\Omega,G_j)\cdot \deg(g,G_j,y)
    \end{align*}
    \end{description}
\end{proof}

Eine wichtige Anwendung der Produktformel ist das folgende
\begin{theorem}[Jordanische Kurvensatz]\label{2.16}
Es seien $C_1$ und $C_2$ zwei kompakte, zueinander Homeomorphe Teilmengen des $\R^n$. Dann besitzen
$\R^n\setminus C_1$ und $\R^n\setminus C_2$ die selbe Anzahl von Zusammenhangskomponenten.
\end{theorem}
 %% Nach dem Einschub geht es weiter mit 2.4 Eigenschaften des Abbildungsgrades.

\section{Der Fixpunktsatz von Kakutani und eine Anwendung in der Spieltheorie}

Im Folgenden beweisen wir die sogenannten \textit{Nash-Gleichgewichte} in n-Personen.
Als Vorbereitung ist es notwendig, den \textit{Brouwer'schen Fixpunktsatz} auf mengenwertige Funktionen
zu verallgemeinern.

Für $K\subset \R^n$ konvex, kompakt, bezeichnen wir mit $\CS(K)$ die Menge der konvexen, nichtleeren
(Teil-)Mengen von $K$.

\begin{theorem}[Kakutani]\label{2.16}
    Sei $K\subset \R^n$ kompakt, konvex $f:K\ra \CS(K)$. Falls die Menge
    \[
        \Gamma := \{(x,y)\in K\times K\, |\, y\in f(x)\}
    \]
    in $K\times K$ abgeschlossen ist, dann existiert $x\in K$ mit $x\in f(x)$.
\end{theorem}

\begin{proof}
    In einer Dimension: $K=[v_1,v_2]$, wir wählen $y_j\in f(v_j)$ und definieren
    \[
        f^1(x)=\sum_{j=1}^{2} \lambda_j(x)y_j,
    \]
    wobei $\lambda _j(x)$ die baryzentrischen Koordinaten von $x$ bezeichne, d.h.
    \[
        \lambda_j(x)\geq 0, \qq \sum_{i=1}^2 \lambda_1(x)=1,\qq x=\sum_{j=1}^2 \lambda_j(x)v_j
    \]
    Nach Konstruktion ist $f^1:K\ra K$ stetig, besitzt mit Brouwer also einen Fixpunkt. Das hilf
    noch nicht viel. Wir betrachten die $k-1$-te baryzentrische Unterteilung der Menge $K$ gegeben
    durch die Vertizes 
    \[
        \{v_i\}_{i=1}^{2^{k-1}+1}
    \]
    Wir wählen wieder $y_j\in f(v_j)$ und definieren $f^k(v_j)=g_j$ interpolieren auf den unterteilten
    Intervallen wie gehabt durch baryzentrische Koordintaten. Es gilt
    \[
        f^k: K\ra K
    \] 
    ist stetig, also mit Fixpunkt
    \begin{align}\label{7}
        x^k=\sum_{i=1}^2 \lambda_i^kv_i^k= \sum_{i=1}^2 \lambda_i^ky_i^k \qq \text{mit} \qq 
        y_i^k=f^k(v_i^k)
    \end{align}
    in einem Teilintervall  (in einem Simplex der $k$-fachen baryzentrischen Koordinaten). Es gilt:
    \[
        \left(x^k,\lambda_1^k,\lambda_2^k,y_1^k,y_2^2\right)\in K\times[0,1]^2\times K^2
    \]
    Somit konvergiert eine Teilfolge in $K$ gegen
    \[
        \left(x^0,\lambda_1^0,\lambda_2^0,y_1^0,y_2^0\right)
    \]
    Nachdem die Teilintervalle zu Punkten degenerieren gilt:
    \[
        v_i^k\ra x^0 \qq i=1,2
    \]
    Wir haben 
    \[
        \left(v_i^k,y_i^k\right)\in\Gamma\ra \left(v_i^0,y_i^0\right)\in \Gamma \qq 
        \text{wegen Abgeschlossenheit von }\Gamma
    \]
    Somit gilt $y_i^0\in f(x^0)$, $i=1,2$. Mit (\ref{7}) folgt
    \[
        x_0 =\lim_{k\ra\infty} \sum_{i=1}^2\lambda_i^ky_i^k=\sum_{i=1}^2 \lambda^0_iy_i^0 
        \in f\left(x^0\right),
    \]
    $f\left(x^0\right)$ ist konvex. 

    In $n$ Dimensionen nehmen wir zunächst an, dass $K$ ein Simplex (Konvexe Hülle von $n+1$
    nichtdegenerierten Vertizes) ist. Wir wenden wiederholt die Baryzentrische Unterteilung an,
    der Durchmesser der Simplizes schrumpft mit einem Faktor $\frac{n}{n+1}$.
    Zur $k$-ten baryzentrischon Unterteilung mit Vertizes $\{v_j\}_{j=1}^{N_k}$. Wählen wir wieder
    $y_j\in f(v_j)$ mit $f^k(x)$ durch Interpolation mittels baryzentrischen Koordinaten in jdn.
    Unter-Simplex. Wir erhalten wieder Fixpunkte $x^k$, und eine Teilfolge von
    \[
        (x^k,\lambda_1^k,…,\lambda_{n+1}^k,y_1^k,…,y_{n+1}^k)
    \]
    konvergiert im „richtigen“ Unter-Simplex. Der Rest des Beweises folgt exakt wie im eindimensionalen
    Fall. Ist $K$ kein Simplex, so können wir $f^k$ jeweils auf einen Simplex $\tilde K\supset K$ stetig
    fortsetzen und verfahren wie im Beweis des \textit{Brouwer'schen Fixpunktsatzes}.
    \[ \]
\end{proof}

\subsection*{Spieltheorie:}

Ein $n$-Personen Spiel besteht aus $n$ Spielern, wobei der $i$-te Spieler $m_i$ mögliche Aktionen
ausführen kann. Die Menge der möglichen Aktionen des $i$-ten Spielers bezeichnen wir mit
$\Phi_i=\{ 1,…,m_i \}$. Nachdem jeder Spieler eine Aktion ausgeführt hat, wird abgerechnet. Der $i$-te
Spiele bekommt den Payoff
\[
    R_i(\vp), \qq \vp=(\vp_1,…,\vp_n)\in \Phi=\prod_{j=1}^n \Phi_j
\]
Wir betrachten den Fall, dass das Spiel häufig ausgeführt wird, und jeder Spieler seine Aktion nach
einer Wahrscheinlichkeitsverteilung (einer Strategie) auf $\Phi_i $ wählt. Diese bezeichenen wir mit
\[
    S_i=\{s_i^1,…,s_i^{m_i}\},\qq s_i^k\geq 0, \qq \sum_{k=1}^{m_i}s_i^k=1 \qq \forall i
\]
Die Menge aller möglichen Strategien für den $i$-ten Spieler nennen wir $S_i$. $S_i^k$ ist die
Wahrscheinlichkeit, dass der $i$-te Spieler die Aktion $k$ wählt. Die Wahrscheinlichkeit, dass jeder
Spieler eine bestimmte Aktion $s_i^k$ wählt ist
\[
    S(\vp)=\prod_{i=1}^n s_i(\vp), \qq s_i(\vp)=s_i^{k_i}
\]
mit $\vp=(k_1,…,k_n)\in \Phi$. Wir nehmen hier an, dass die Spieler ihre Aktionen unabhängig voneinander
wählen. Der mittlere Payoff für Spieler $i$ ist dann der Erwartungswert für $R_i(\vp)$. Wir nennen diesen
$R_i(s)$ und erhalten
\[
    R_i(s)=\sum_{\vp\in\Phi}s(\vp) K_i(\vp)
\]
Nach Konstruktion ist
\[
    R_i(s)
\]
stetig in $s$. Was ist nun die optimale Strategie für Spiel $i$? Falls alle anderen Spiele nach einer
bekannten  Strategie handeln, ist es optimal $\ol{s}_i$ so zu wählen, dass
\begin{align}\label{8}
    R_i(s\setminus \ol{s_i})=\max_{\tilde s_i \in S_i} R_i(s\setminus \tilde s_i)
\end{align}
$s\setminus \tilde s_i$: Strategiekomination, die sich ergibt, wenn man $s_i$ durch  $\tilde s_i$
ersetzt. Wir bezeichnen mit $B_i(s)$ die Menge aller Strategien $\ol s_i\in S_i$, die $(\ref{8})$ erfüllen.
Insbesondere gilt $\ol s_i\in B_i(s)$ genau dann wenn $\ol s_i^k=0$ falls
\[
    R_i(s\setminus (0,…,1,…)) < \max_{1<l<m_i}R_i(s\setminus(0,…,1,…)) \qq \text{mit der 1 als $l$-ter
    Eintrag}
\]
Insbesondere existiert immer eine reine (immer die selbe Aktion wird gewählt) optimale Strategie
zu gegebenem $S$. Weiter sieht man, dass $B_i(s)$ konvex ist und nicht leer. Es seien nun $s,\ol s\in S$.
Wir nennen $\ol s$ die optimale Antwort gegen $s$, falls $\ol s_i\in B_1(s)$ $i=1,…,n$.
Es sei $B(s)$ die Menge aller optimalen Antworten gegen $s$, wir haben
\[
    B(s)=\prod_{i=1}^n B_i(s).
\]

\begin{defi}[Nash-Gleichgewicht]
    Eine Strategiekombination $\ol s\in S$ heißt Nash-Gleichgewicht, falls gilt
    \[
        \ol s \in B(\ol s)
    \]
    d.h. $\ol s$ ist eine optimale Antwort gegen sich selbst. Anders gesagt, $\ol s$ ist ein 
    Nash-Gleichgewicht, falls kein Spieler durch eigene Änderung seiner Strategie seinen Payoff erhöhen
    kann.
\end{defi}

\subsection*{Beispiele:}

%\begin{description}
    \textbf{Das Gefangenendilemma:}
    Zwei Gefangene können entweder schweigen, oder mit der Polizei kooperieren, das heißt den jeweils
    anderen belasten.
    \[
        \begin{array}{c|cc||c|cc} R_1& k_2& s_2& R_2&k_2&s_2\\ \hline k_1&0&2&k_1&0&-1\\
            s_1&-1&1&s_1&2&1\end{array}
    \]
%\end{description}
\begin{theorem}\label{2.18}
    Jedes $n$-Personen Spiel besitzt ein Nash-Gleichgewicht.
\end{theorem}

\begin{proof}
    Wir betrachten $s\mapsto B(s)$, $S$ ist kompakt, konvex, $B(s)\subset S$ ist nicht leer, konvex für
    alle $s$. Abgeschlossenheit des Graphen folgt mit Folgenkompaktheit.
    \[
        s^n\in S, \qq \ol s^m\in B(s^m), \qq s^m\ra s, \qq \ol s^k\ra \ol s
    \]
    Wir haben $R_i(s^m\setminus \tilde s_i)\leq R_i(s^m\setminus \ol s_m)$ für alle $\tilde s_i \in S_i$
    mit Stetigkeit von $R_i$ folgt sofort
    \[
        R_i(s\setminus \tilde s_i) \leq R_i(s\setminus \ol s) \qq \forall s_i\in S_i
    \]
    Somit ist $\ol s \in B(s)$ und $\Gamma= \{ (s,B(s))\subset S^2 \}$ ist abgeschlossen.
    Der \textit{Fixpunktsatz von Kakutani} garantiert Existenz von
    \[
        \ol s\in S \qq \text{mit} \qq \ol s \in B(\ol s).
    \]
\end{proof}

\section{Der Leray-Schauder Grad und der Schauder'sche Fixpunktsatz}

Das Ziel ist es, den Fixpunktsatz von Brouwer auf $\infty$-dimensionale Räume zu verallgemeinern.
So einfach ist das nicht.

\begin{theorem}[Kakutani]\label{3.1}
    Sei $H$ ein $\infty$-dimensionaler separabler Hilbertraum. Dann existiert eine stetige Abbildung
    \[
        f:H\ra H,
    \]
    welche die Einheitenkugel in sich selbst abbildet und die keinen Fixpunkt besitzt.
\end{theorem}

\begin{proof}
    Sei $(e_n)_{n\in \Z}$ eine Orhonormalbasis von $H$. Für $x=\sum_{j=-\infty}^\infty\alpha_je_j$
    definieren wir
    \[
        U(x):=\sum_{j=-\infty}^\infty \alpha _je_{j+1};\qq f(x):=\frac12 (1-\|x\|)e_0+U(x)
    \]
    Zu zeigen:
    \begin{description}
    \item{-}
    $U$ ist stetig.
    \item{-}
    $\|U(x)\|<1$ für $\|x\|<1$
    \item{-}
    $f$ besitzt keinen Fixpunkt in $\ol{B_1(0)}$
    \begin{itemize}
        \item bei $x=0$.
        \item auf dem Rand.
        \item im Innern.
    \end{itemize}
    \end{description}
    \[  \]
\end{proof}

\subsection{Der Abbildungsgrad auf endlichen Banachräumen}

\begin{defi}\label{3.2}
    Sei $X$ ein (reeller) Banachraum von Dimension $n$, und sei $\Phi$ ein Isomorphismus von $X$ nach 
    $\R^n$. Für $f\in D_y(\ol \Omega,X)$, $\Omega \subset X$ offen, beschränkt, $y\in X$ definieren wir
    \[
        \deg(f,\Omega,y)=\deg(\Phi\circ f \circ \Phi^{-1},\Phi(\Omega),\Phi(y))
    \]
\end{defi}

\begin{prop}\label{3.3}
    Der in Definition \ref{3.2} definierte Abbildungsgrad ist unabhängig von der Wahl von $\Phi$.
\end{prop}

\begin{proof}
    Sei $\Psi$ ein zweiter Isomorphismus. Es gilt $A=\Psi\circ\Phi^{-1}\in \Gl(n)$. Wir schreiben
    \[
        f^\star:=\Phi \circ f \circ \Phi^{-1}, \qq y^\star\in \RV(\tilde f^\star)
    \]
    und wählen wie üblich $\tilde f^\star\in D^1(\Phi(\ol\Omega),\R^n)$ aus derselben Komponente von
    $D_y(\Phi(\ol \Omega),\R^n)$ wie $f^\star$, so dass $y^\star\in \RV(\tilde f^\star)$. Dann ist auch
    \[
        A\circ \tilde f^\star \circ A^{-1}\in D^1(\Psi (\ol \Omega),\R^n)
    \]
    in derselben Komponente von $D_y(\Psi(\ol \Omega),\R^n)$ wie $A\circ f^\star\circ A^{-1}=\Psi\circ
    f\circ \Psi^{-1}$.
    Nach der Kettenregel gilt
    \begin{align*}
        J_{A\circ \tilde f^\star\circ A^{-1}}(Ay^\star)&= \det(A) J_{\tilde f^\star}(y^\star)
        \det(A^{-1})=J_{\tilde f^\star} (y^\star).
    \end{align*}
    Und es folgt
    \[ 
        \deg(\Psi\circ f \circ\Psi^{-1},\Psi(\Omega),\Psi(y))
        = \deg(\Phi\circ f \circ\Phi^{-1},\Phi(\Omega),\Phi(y))
    \]
\end{proof}

Die Eigenschaften des Abbildungsgrades im $\R^n$ übertragen sich mit Definition \ref{3.2} auf den
Abbildungsgrad in endlichdimensionalen Banachräumen. Das gilt ebenfalls für die Reduktionseigenschaft aus
Satz \ref{2.14}, nachdem $\Phi:X\ra\R^n$ so gewählt werden kann, dass $\Phi(X_1)=\R^n$, wobei $X_1$ ein
$m$-dimensionaler Untervektorraum von $X$ ist.

\subsection{Kompakte Operatoren}

Es seien $X$ und $Y$ Banachräume, $\Omega\subset X$ (nicht unbedingt offen).

\begin{defi}
    Ein stetiger Operator $F:\Omega\ra Y$ heißt
    \begin{description}
        \item{•} \textit{endlichdimenisional}, falls ein endlichdimensionaler Untervektorraum $Y_1\subset Y$
        existiert, sodass $F(\Omega)\subset Y_1$.
        \item{•} \textit{kompakt}, falls $F$ beschränkte Teilmenge von $\Omega$ auf relativ kompakte
        Teilmengen von $Y$ abbildet.
    \end{description}
    Wir bezeichnen die Menge der endlichdimensionalen Operatoren von $\Omega$ nach $Y$ mit
    \[
        \ms F(\Omega,Y)
    \]
    und die kompakten mit
    \[
        \ms C(\Omega,Y)
    \]
\end{defi}

\begin{remark}
    \begin{description}
        \item{1)}
        Es gilt
        \[
            \ms F(\Omega,Y)\subset \ms C(\Omega,Y)\subset C(\Omega,Y)
        \]
        \item{2)}
        Falls $\Omega$ kompakt ist, so gilt
        \[
            \ms C(\Omega,Y)\subset C(\Omega,Y)
        \]
        \item{3)}
        Falls $\dim(Y)<\infty$, so gilt
        \[
            \ms F(\Omega,Y)=\ms (\Omega,Y)
        \]
        \item{4)}
        Falls $\Omega\subset\R^n$ beschränkt, so gilt
        \[
            \ms F (\ol \Omega,\R^n)=\ms C(\ol \Omega,\R^n)=C(\ol\Omega,\R^n)
        \]
    \end{description}
\end{remark}

\subsubsection*{Korrekturen}

Bei Satz \ref{2.15}. Vorsicht vor der unbeschränkten Komponente $K_\infty$ von $\R^n\setminus f(\partial
\Omega)$! Diese taucht jedoch in der Produktformel ausschließend nicht auf, denn
\[
    f(\ol \Omega)\subset B_r(0), \, \text{mit}\, f(\ol \Omega) \text{ kompakt}
\]
\[
    \Ra \, \deg(f,\Omega,K_\infty\cap B_{2r}(0))=0 \qq \text{und} \qq f(\ol \Omega) \cap (\R^n\setminus
    B_{2r}(0))=\varnothing
\]
somit können wir vereinfacht schreiben
\[
    \deg(f,\Omega,K_\infty)=0
\]
Der in der Produktformel aufauchende Term
\[
    \deg(f,\Omega,K_\infty)\cdot \deg(g,K_\infty,y)
\]
ist also unproblematisch, da er verschwindet.\\[0.5cm]

\begin{lem}\label{3.5}
    Sei $K\subset X$ kompakt, $\eps>0$. Dann existiert ein endlichdimensionaler Untervektorraum $X_\eps
    \subset X$ und eine stetige Abbildung
    \[
        P_\eps:K\ra X_\eps,
    \]
    so dass
    \[
        \|P_\eps(x)-x\|\leq \eps \qq \forall x\in K
    \]
\end{lem}

\begin{proof}
    Es sei $\{x_j\}_{j=1}^N\subset K$, so dass $\bigcap _{j=1}^NB_\eps(x_j)$ eine Überdeckung von $K$
    darstellt. Sei $\{\Phi_j\}_{j=1}^N$ eine Zerleung der Eins, welche $\{B_\eps(x_j)\}_{j=1}^N$
    untergeordnet ist. Wir definieren
    \[
        P_\eps(x):= \sum_{j=1}^N \Phi_j(x)x_j
    \]
    Damit gilt
    \[
        \|P_\eps(x)-x\|=\|\sum_{j=1}^N\Phi_j(x)x_j-\sum_{j=1}^N \Phi_j(x)x_j\leq \sum_{j=1}^N \Phi_j(x)
        \|x-x_j\|\leq \eps\|
    \]
\end{proof}

\begin{theorem}\label{3.6}
    $X,Y$ Banachräume. Sei $\Omega\subset X$ beschränkt. Dann gilt
    \[
        \ol{\ms F(\Omega,Y)}^{C{\Omega,Y}}=\ms C(\Omega,Y)
    \]
    (D.h. der Abschluss von $\ms F(\Omega,Y)$ bzgl. der Norm der gleichmäßigen Konvergenz ist der Raum
     der kompakten Operatoren $\ms C(\Omega,Y)$.)
\end{theorem}

\begin{proof}
    Es sei $(F_n)_{n\in \N}$ eine Folge in $\ms C(\Omega,Y)$, sodass  $F_n\ra F$. Wir nehmen an, dass
    $F\nin \ms C(\Omega,Y)$. Dann existiert eine Folge $(x_n)_{n\in \N}$, so dass
    \[
        \|F(x_n)-F(x_l)\|\geq \rho>0 \qq\text{für}\qq k\neq l
    \]
    Es sei $n$ so groß, dass
    \[
        \|F_n-F\|_{\ms C(\Omega,Y)}<\frac{\rho}{4}
    \]
    Damit gilt
    \[
        \|F_n(x_k)-F_n(x_l)\|\geq \frac\rho2.
    \]
    Das ist ein Widerspruch zur Annahme, dass $F_n\in \ms C(\Omega,Y)$. Somit gilt auch
    \[
        \ol{\ms F(\Omega,Y)}^{\ms C(\Omega,Y)}\subset \ms C(\Omega,Y)
    \]
    Sei andererseits $F\in \ms C(\Omega,Y)$,und sei $K=\ol{F(\Omega)}$ kompakt. sei $P_\eps$ gewählt wie
    in Lemma \ref{3.5}. Es gilt
    \[
        F_\eps=P_\eps\circ F\in \ms F(\Omega,Y) \qq \text{und}\qq F_\eps\ra F
    \]
\end{proof}

Zur Definition des Abbildungsgrades auf unendlichdimensionalen Räumen betrachten wir insbesondere
kompakte Störungen der Identität (d.h. Operatoren der Form $\Id+F$, $F\in \ms C$). Wir bezeichnen deshalb
zunächst einige Interessante Eigenschaften von $\Id+F$.

\begin{lem}\label{3.7}
    Sei $X$ ein Banachraum und $\Omega\subset X$ beschränkt und abgeschlossen, $F\in \ms C(\Omega,Y)$.
    Dann ist $\Id + F$ eine eigentliche Abbildung (d.h. $(\Id+F)^{-1}(K)$ ist kompakt für $K\in X$ 
    kompakt), die abgeschlossene Teilmengen von $\Omega$ auf abgeschlossene Teilmengen von $X$ abbildet.
\end{lem}

\begin{proof}
    Sei $A\subset \Omega$ abgeschlossen, und sei
    \[
        y_n=(\Id+F)(x_n), \qq \text{mit } x_n \in A,
    \]
    so dass $y_n \ra y$ in $X$. Zu zeigen ist, dass $y\in (\Id+F)(A)$. Es gilt
    \[
        y_n-x_n=F(x_n),
    \]
    somit folgt $y_n-x_n\ra z$ in $X$, nach Extraktion einer Teilfolge. Es folgt $x_n\ra x \in A$, da
    $A$ abgeschlossen und $y_n$ konvergent nach Annahme. Wir haben aber $x=y-z\in A$. Nachdem aber
    $y=x+F(x)\in (Id+F)(A)$ folgt die Behauptung.

    Sei nun $\Omega$ abgeschlossen und $K\subset X$ kompakt. Sei $(x_n)_{n\in \N}$ eine Folge in
    $(\Id+F)^{-1}(K)$. Wir können wieder eine Teilfolge
    \[
        y_n=x_n+F(x_n)
    \]
    wählen, so dass $y_n\ra y$. Wie im vorhergehenden Teil folgt, folgt dass $x_n\ra x$, somit ist
    $(\Id+F)^{-1}(K)$ kompakt.
    \[ \]
\end{proof}

\subsection{Der Leray-Schauder Grad}

Sei $\Omega\subset X$, $X$ ein Banachraum. Wir setzten
\[
    \ms D_y(\ol \Omega,X)=\{ F\in \ms C(\ol\Omega,X)\, | \, y\nin(\Id + F)(\partial \Omega) \}
\]
und
\[
    \ms G_y(\ol\Omega,X):=\{F\in \ms F(\ol \Omega,X)\, |\, y\nin(\Id+F)(\partial\Omega)\}
\]
Es gilt für $F\in \ms D_y(\ol \Omega,X)$, dass $\dist(y,(\Id+F)(\partial\Omega))>0$, da $\Id+F$
abgeschlossene Mengen auf abgeschlossene Mengen abbildet.

\begin{defi}{\label{3.8}}
    Sei $\Omega\subset X$ offen, beschränkt, $y\in X$, $F\in \ms D_y(\ol \Omega,X)$. Sei $\rho:=
    \dist(y,(\Id+F)(\partial \Omega))$ und sei
    \[
        F_1\in \ms F(\ol \Omega, X),
    \]
    sodass
    \[
        \|F-F_1\|<\rho. \qq (\Ra \, F_1\in \ms G_y(\ol \Omega,X))
    \]
    Nun sei $X_1\subset X$ ein endlichdimensionaler Untervektorraum, so dass
    \[
        F_1(\Omega)\subset X_1,\, y\in X_1
    \] 
    und sei $\Omega_1=\Omega\cap X_1$. Damit ist $F_1\in \ms G_y(\ol\Omega_1,X_1)$ und wir definieren
    \[
        \deg(\Id+F,\Omega,y):=\deg(\Id+F_1,\Omega_1,y)
    \]
\end{defi}

\begin{prop}\label{3.9}
    die Obige Definiton ist unabhängig von der Wahl von $F_1$ und $X_1$.
\end{prop}

\begin{proof}
    Sei $F_2\in \ms F (\ol \Omega,X)$, $\|F_2-F\|<S$, $X_2$ entsprechend endlichdim UVR von $X$. Es sei
    $X_0=X_1+X_2$, $\Omega_0=\Omega\cap X_0$. Dann gilt
    \[
        F_i\in \ms F_y(\ol\Omega_0,X_0) \qq i=1,2
    \]
    und es folgt wegen der Reduktionseigenschaft des endlich dimenisionalen Abbildungsgrades, dass
    \[
        \deg(\Id + F_i,\Omega_0,y)=\deg(\Id+F_i,\Omega_i,y)\qq i=1,2
    \]
    Sei $H(t)=\Id + (1-t)F_1+F_2.$ Es folgt, dass
    \[
        H(t)\in D_y(\ol\Omega_0,X_0) \qq t\in [0,1],
    \]
    nachdem gilt $\|H(t)-(\Id+F)\|<\rho$ $\forall t\in [0,1]$.
    \[
        \Ra \, \deg(\Id+F_1,\Omega_0,y)=\deg(\Id+F_2,\Omega_0,y)
    \]
    Somit ergibt Definiton \ref{3.8} Sinn.
    \[ \]
\end{proof}

\begin{theorem}\label{3.10}
    Sei $\Omega\subset X$ beschränkt und offen und sei $F\in \ms D_y(\ol\Omega,X)$, $y\in X$. Dann gilt
    \begin{description}
        \item{i)}
        $\deg(\Id+F,\Omega,y)=\deg(\Id+F-y,\Omega,0)$
        \item{ii)}
        $\deg(\Id, \Omega,y)=1$, falls $y\in \Omega.$
        \item{iii)}
        Falls $\Omega_1,\Omega_2$ offene, disjunkte Teilmengen von $X$ sind, so dass
        $y\nin (\Id+F)(\ol\Omega\setminus (\Omega_1\cup\Omega_2)),$ dann gilt
        \[
            \deg(\Id+F,\Omega,y)=\deg(\Id+F,\Omega_1,y)+\deg(\Id+F,\Omega_2,y).
        \]
        \item{iv)}
        Falls $H:[0,1]\times\ol\Omega\ra X$ und $y:[0,1]\ra X$ beide stetig, so dass $H(t)\in 
        \ms D_{y(t)}(\ol\Omega,X) \qq \forall t\in [0,1]$. Dann gilt
        \[
            \deg(\Id+H(0),\Omega,y(0))=\deg(\Id+H(1), \Omega,y(1))
        \]
    \end{description}
\end{theorem}

\begin{proof}
    \textit{i)-iii)} folgen sofort aus der Definition und den entsprechenden Eigenschaften des
    endlichdimensionalen Abbildungsgrades.

    \noindent Zu \textit{iv)}: Mittels \textit{i)} betrachten wir $y(t)=0$ ($\tilde H(t)=H(t)-y(x)$).
    Nachdem $H([0,1],\partial \Omega)$ abgeschlossen und kompakt ist, gilt
    $\rho=\dist(y,H([0,1])\times\partial\Omega)>0$. Nach Satz \ref{3.6} existiert $H_1\in \ms F([0,1]
    \times\ol\Omega, X)$, so dass $\|H(t)-H_1(t)\|<\rho$ $t\in [0,1]$. Damit folgt
    \[
        \deg(\Id+H(t),\Omega,0)=\deg(\Id+H_1(t),\Omega,0) \qq \forall t\in [0,1].
    \]
    Die Behauptung folgt nun wieder aus den Eigenschaften des endlichen Abbildungsgrades.
    \[ \]
\end{proof}


\end{document}
